\documentclass{article}
\usepackage{iclr2027_conference,times}
\usepackage[T1]{fontenc}
\input{math_commands.tex}
\usepackage{hyperref}
\hypersetup{pdfauthor={},pdfcreator={},pdfproducer={},pdfsubject={},pdfkeywords={}}
% anonymity: no creation/modification date, trailer ID or TeX banner in the PDF
\ifdefined\pdfinfoomitdate \pdfinfoomitdate=1 \fi
\ifdefined\pdftrailerid \pdftrailerid{} \fi
\ifdefined\pdfsuppressptexinfo \pdfsuppressptexinfo=-1 \fi
\usepackage{url}
\usepackage{amsmath,amssymb,amsthm}
\usepackage{booktabs}
\usepackage{graphicx}
\usepackage{xcolor}
\usepackage{enumitem}
\usepackage{multirow}
\usepackage{longtable}

% ---------------------------------------------------------------- numbers
% Every number in the text is a macro. numbers.tex is generated from the result files by exp/make_numbers.py
% (\newcommand, read first, wins); numbers_placeholder.tex (\providecommand only) fills whatever is still
% undefined. Rebuild everything with: python paper/exp/build_all.py
\InputIfFileExists{numbers.tex}{}{}
\InputIfFileExists{numbers_story.tex}{}{}
\InputIfFileExists{numbers_coverage.tex}{}{}
% numbers_placeholder.tex -- fallback values for every number in main.tex.
%
% Load order in main.tex:  numbers.tex (if it exists; generated by exp/make_numbers.py with \newcommand)
%                          -> this file.
% Everything here is \providecommand, so a macro defined by an earlier file wins and this file only fills
% gaps. Values below are the prior-grid values (results/THINKING_VS_DATA.md, results/GATE_FINETUNE_STUDY.md,
% results/LOGIC_FINETUNE.md) or \TBD for cells preregistered in PREDICTIONS.md and not yet run.
% Percentages are written without the % sign; the text adds \%.

\providecommand{\TBD}{\textbf{[TBD]}}

% ---------------------------------------------------------------- setup
\providecommand{\nTrain}{180}
\providecommand{\nTrainSmall}{60}
\providecommand{\nTrainLarge}{540}
\providecommand{\nTest}{240}
\providecommand{\nTestClass}{120}
\providecommand{\loraRank}{8}
\providecommand{\loraAlpha}{16}
\providecommand{\loraDropout}{0.05}
\providecommand{\nEpochs}{3}
\providecommand{\learningRate}{2\times 10^{-4}}
\providecommand{\ciHalf}{6}
\providecommand{\nShots}{4}
\providecommand{\nForms}{12}
\providecommand{\nFormsHalf}{6}
\providecommand{\maxNewPrime}{512}
\providecommand{\maxNewDiv}{176}
\providecommand{\maxNewValid}{320}
\providecommand{\nPriorRuns}{36}
\providecommand{\priorWallMin}{319}
\providecommand{\priorVerdictDate}{2026-08-04}
\providecommand{\alphaLevel}{0.05}

% ---------------------------------------------------------------- trace structure and training tokens
\providecommand{\divsevenSteps}{4}
\providecommand{\divsevenSixSteps}{6}
\providecommand{\divsevenTableSize}{70}
\providecommand{\divsevenStepSamples}{720}
\providecommand{\primeStepsMin}{11}
\providecommand{\primeStepsMax}{25}
\providecommand{\primeStepsMeanPrimes}{19.2}
\providecommand{\primeStepsMedComp}{1}
\providecommand{\primeStepsMeanComp}{2.6}
\providecommand{\primeTraceTok}{71.5}
\providecommand{\divsevenTraceTok}{58.0}
\providecommand{\validTraceTok}{37.0}
\providecommand{\tokA}{360}
\providecommand{\primeTokB}{13{,}300}
\providecommand{\divsevenTokB}{10{,}800}
\providecommand{\validTokB}{7{,}000}
\providecommand{\tokRatioLow}{19}
\providecommand{\tokRatioHigh}{37}
\providecommand{\runMinLow}{2}
\providecommand{\runMinHigh}{17}
\providecommand{\validPool}{1{,}600}
\providecommand{\cTokDiff}{6}
\providecommand{\baseEvalRatio}{38}
\providecommand{\nProbeFeatures}{40}

% ---------------------------------------------------------------- constants of the tasks (exact, not results)
\providecommand{\condLow}{0.10}
\providecommand{\condHigh}{0.22}
\providecommand{\condMarg}{0.14}
\providecommand{\hardNegPop}{996}
\providecommand{\compPop}{7{,}939}
\providecommand{\hardNegFrac}{12.5}
\providecommand{\winMargin}{10}
\providecommand{\neitherMargin}{8}

% ---------------------------------------------------------------- surface probe and rules (prime test set)
% surface probes: experiments/analysis_probe.py (results/analysis/probe.json). prime values are its seed-0
% output (best by 5-fold CV = logreg on one-hot + engineered features; one-hot logreg alone). div7 and valid
% are stand-ins from a plain one-hot / atom-renamed n-gram logreg (C = 1) until analysis_probe.py covers them.
\providecommand{\primeProbe}{89.2}
\providecommand{\primeProbeOneHot}{78.8}
\providecommand{\divsevenProbe}{50.0}
\providecommand{\validProbe}{100.0}
% rules evaluated on this paper's prime test set (where A = 82.3)
\providecommand{\ruleOdd}{76.2}
\providecommand{\ruleNoTwoThree}{86.7}
\providecommand{\ruleNoSeven}{93.8}
% the same rules on the earlier answer-only study's prime test set (where answers gave 82.9)
\providecommand{\ruleOddPrior}{78.8}
\providecommand{\ruleNoTwoThreePrior}{89.2}
\providecommand{\primeAagreeOdd}{93.5}
\providecommand{\ruleLastDigit}{82.5}
\providecommand{\primeAagreeLastDigitLo}{99.6}
\providecommand{\primeAagreeLastDigitHi}{100.0}
\providecommand{\primeAprimesCorrect}{99.6}
\providecommand{\primeAcompCorrect}{65.0}
\providecommand{\primeHardNegN}{15}
\providecommand{\primeHardNegAcorrect}{0 and 0}
\providecommand{\primeBprimesCorrect}{19.6}
\providecommand{\primeBcompCorrect}{95.8}
\providecommand{\primeBlargePrimes}{46.7}
\providecommand{\primeBlargeComp}{90.8}

% ---------------------------------------------------------------- prime (Qwen2.5-1.5B, n = 180 unless noted)
\providecommand{\primeBase}{59.6}
\providecommand{\primeA}{82.3}
\providecommand{\primeAsZero}{82.1}
\providecommand{\primeAsOne}{82.5}
\providecommand{\primeB}{57.7}
\providecommand{\primeBsZero}{57.9}
\providecommand{\primeBsOne}{57.5}
\providecommand{\primeC}{52.1}
\providecommand{\primeCsZero}{50.0}
\providecommand{\primeCsOne}{54.2}
\providecommand{\primeAsmall}{59.6}
\providecommand{\primeAlarge}{75.0}
\providecommand{\primeBsmall}{57.1}
\providecommand{\primeBlarge}{68.8}
\providecommand{\primeAtok}{80.8}
\providecommand{\primeAtokN}{1{,}320}
\providecommand{\primeBorn}{61.0}
\providecommand{\primeBornN}{9}
\providecommand{\primeBApsZero}{6\times 10^{-7}}
\providecommand{\primeBApsOne}{2\times 10^{-7}}
\providecommand{\primeBCpsZero}{0.22}
\providecommand{\primeBCpsOne}{0.33}

% ---------------------------------------------------------------- div7 (4-digit)
\providecommand{\divsevenBase}{62.5}
\providecommand{\divsevenBaseYes}{84.2}
\providecommand{\divsevenBaseNo}{40.8}
\providecommand{\divsevenA}{51.5}
\providecommand{\divsevenAsZero}{52.1}
\providecommand{\divsevenAsOne}{50.8}
\providecommand{\divsevenAyessZero}{25.8}
\providecommand{\divsevenAnosZero}{78.3}
\providecommand{\divsevenAyessOne}{82.5}
\providecommand{\divsevenAnosOne}{19.2}
\providecommand{\divsevenB}{92.5}
\providecommand{\divsevenBsZero}{92.5}
\providecommand{\divsevenByes}{89.2}
\providecommand{\divsevenBno}{95.8}
\providecommand{\divsevenC}{50.0}
\providecommand{\divsevenCsZero}{49.6}
\providecommand{\divsevenCsOne}{50.4}
\providecommand{\divsevenAsmall}{50.0}
\providecommand{\divsevenAlarge}{50.0}
\providecommand{\divsevenBsmall}{50.0}
\providecommand{\divsevenBApsZero}{4\times 10^{-27}}
\providecommand{\divsevenBApsZeroDisc}{99/2}
\providecommand{\divsevenBCpsZero}{1\times 10^{-25}}
\providecommand{\divsevenABasepsZero}{0.047}
\providecommand{\divsevenABasepsOne}{0.0046}

% ---------------------------------------------------------------- valid (propositional validity, Z3 labels)
\providecommand{\validBase}{30.0}
\providecommand{\validA}{99.6}
\providecommand{\validAsZero}{99.2}
\providecommand{\validAsOne}{100.0}
\providecommand{\validB}{90.0}
\providecommand{\validBsZero}{100.0}
\providecommand{\validBsOne}{80.0}
\providecommand{\validC}{52.5}
\providecommand{\validCsZero}{54.2}
\providecommand{\validCsOne}{50.8}
\providecommand{\validAsmall}{97.5}
\providecommand{\validBsmall}{87.5}
\providecommand{\validAlarge}{100.0}
\providecommand{\validBlarge}{100.0}
\providecommand{\validAtok}{100.0}
\providecommand{\validBApsZero}{0.50}
\providecommand{\validBApsOne}{7\times 10^{-15}}
\providecommand{\validBCpsZero}{2\times 10^{-33}}
\providecommand{\validBCpsOne}{5\times 10^{-10}}

% ---------------------------------------------------------------- earlier answer-only study (shuffled-label control)
\providecommand{\priorTrainN}{240}
\providecommand{\shufMin}{48.5}
\providecommand{\shufMax}{52.5}
\providecommand{\priorCorrectMax}{99.8}
\providecommand{\priorDivthreeBase}{55.0}
\providecommand{\priorDivthreeA}{67.3}
\providecommand{\priorDivthreeShuf}{49.6}
\providecommand{\priorDivsevenBase}{53.3}
\providecommand{\priorDivsevenA}{51.0}
\providecommand{\priorDivsevenShuf}{52.5}
\providecommand{\priorDivthirteenBase}{54.2}
\providecommand{\priorDivthirteenA}{50.0}
\providecommand{\priorDivthirteenShuf}{50.0}
\providecommand{\priorSquareBase}{51.2}
\providecommand{\priorSquareA}{83.3}
\providecommand{\priorSquareShuf}{51.5}
\providecommand{\priorPrimeBase}{50.0}
\providecommand{\priorPrimeA}{82.9}
\providecommand{\priorPrimeShuf}{50.0}
\providecommand{\priorValidBase}{66.7}
\providecommand{\priorValidA}{99.4}
\providecommand{\priorValidShuf}{50.0}
\providecommand{\priorConsistBase}{50.0}
\providecommand{\priorConsistA}{99.8}
\providecommand{\priorConsistShuf}{48.5}

% ================================================================ preregistered cells (PREDICTIONS.md), not yet run
% S1: div7 replication
\providecommand{\divsevenBsOne}{\TBD}
\providecommand{\divsevenBsTwo}{\TBD}
\providecommand{\divsevenAsTwo}{\TBD}
\providecommand{\divsevenBApsOne}{\TBD}
\providecommand{\divsevenBApsTwo}{\TBD}
\providecommand{\verdictSOne}{\TBD}
% S2/S3: divisor family (4-digit, k = 4); mean of two seeds
\providecommand{\divtwoA}{\TBD}
\providecommand{\divtwoB}{\TBD}
\providecommand{\divtwoProbe}{\TBD}
\providecommand{\divthreeA}{\TBD}
\providecommand{\divthreeB}{\TBD}
\providecommand{\divthreeProbe}{\TBD}
\providecommand{\divelevenA}{\TBD}
\providecommand{\divelevenB}{\TBD}
\providecommand{\divelevenProbe}{\TBD}
\providecommand{\divthirteenA}{\TBD}
\providecommand{\divthirteenB}{\TBD}
\providecommand{\divthirteenProbe}{\TBD}
\providecommand{\verdictSTwo}{\TBD}
\providecommand{\verdictSThree}{\TBD}
% S4: scale
\providecommand{\threeBdivsevenA}{\TBD}
\providecommand{\threeBdivsevenB}{\TBD}
\providecommand{\threeBprimeA}{\TBD}
\providecommand{\threeBprimeB}{\TBD}
\providecommand{\threeBvalidA}{\TBD}
\providecommand{\threeBvalidB}{\TBD}
\providecommand{\halfBdivsevenA}{\TBD}
\providecommand{\halfBdivsevenB}{\TBD}
\providecommand{\sevenBdivsevenA}{\TBD}
\providecommand{\sevenBdivsevenB}{\TBD}
\providecommand{\verdictSFour}{\TBD}
% S5: length test (6-digit div7, k = 6)
\providecommand{\stepAcc}{\TBD}
\providecommand{\divsevenSixPred}{\TBD}
\providecommand{\divsevenSixPredBal}{\TBD}
\providecommand{\divsevenSixA}{\TBD}
\providecommand{\divsevenSixB}{\TBD}
\providecommand{\divsevenSixByes}{\TBD}
\providecommand{\verdictSFive}{\TBD}
% S6: prime hard negatives (odd composites with no prime factor <= 7)
\providecommand{\primeHardN}{\TBD}
\providecommand{\primeHardA}{\TBD}
\providecommand{\primeHardB}{\TBD}
\providecommand{\verdictSSix}{\TBD}
% answer-then-trace (B') and test-time-compute controls on the base model
\providecommand{\divsevenBp}{\TBD}
\providecommand{\primeCot}{\TBD}
\providecommand{\divsevenCot}{\TBD}
\providecommand{\validCot}{\TBD}
\providecommand{\primeFew}{\TBD}
\providecommand{\divsevenFew}{\TBD}
\providecommand{\validFew}{\TBD}

% added with the S1 update (generated by make_numbers.py; placeholders only)
\providecommand{\divsevenBnSeeds}{3}
\providecommand{\divsevenBApMax}{\TBD}
\providecommand{\divsevenBCpsOne}{\TBD}
\providecommand{\divsevenBsZeroRerun}{\TBD}
\providecommand{\stepLookupShare}{\TBD}
\providecommand{\divsevenFourPred}{\TBD}
\providecommand{\divsevenFullTrace}{\TBD}
\providecommand{\divsevenFourPredBal}{\TBD}
\providecommand{\divsevenTransferSixB}{\TBD}
\providecommand{\primeHardAall}{\TBD}
\providecommand{\divsevenBCpMax}{\TBD}
\providecommand{\stepGenSeeds}{\TBD}
\providecommand{\divsevenGenAcc}{\TBD}
\providecommand{\divsevenDigitCopy}{\TBD}
\providecommand{\divsevenCarry}{\TBD}
\providecommand{\divsevenAccGivenWrong}{\TBD}
\providecommand{\primeHardAneg}{\TBD}
\providecommand{\primeHardNneg}{\TBD}
\providecommand{\divsevenBsZero}{92.5}

% added after an internal review pass run with AI agents: fallbacks for macros generated by make_numbers.py
\providecommand{\divelevenBAci}{\TBD}
\providecommand{\divelevenBAdiff}{\TBD}
\providecommand{\divelevenBApsZero}{\TBD}
\providecommand{\divelevenBase}{\TBD}
\providecommand{\divelevenC}{\TBD}
\providecommand{\divsevenABasepsTwo}{\TBD}
\providecommand{\divsevenAsd}{\TBD}
\providecommand{\divsevenBAci}{\TBD}
\providecommand{\divsevenBAdiff}{\TBD}
\providecommand{\divsevenBsd}{\TBD}
\providecommand{\divsevenCsd}{\TBD}
\providecommand{\divsevenD}{\TBD}
\providecommand{\divsevenSelfTrace}{\TBD}
\providecommand{\divsevenSixBAci}{\TBD}
\providecommand{\divsevenSixBAdiff}{\TBD}
\providecommand{\divsevenSixBApsZero}{\TBD}
\providecommand{\divsevenSixBase}{\TBD}
\providecommand{\divsevenSixC}{\TBD}
\providecommand{\divsevenSixPredGuess}{\TBD}
\providecommand{\divsevenSixProbe}{\TBD}
\providecommand{\divsevenSixTraceTok}{\TBD}
\providecommand{\divthirteenAmid}{\TBD}
\providecommand{\divthirteenBAci}{\TBD}
\providecommand{\divthirteenBAdiff}{\TBD}
\providecommand{\divthirteenBApsZero}{\TBD}
\providecommand{\divthirteenBase}{\TBD}
\providecommand{\divthirteenBmid}{\TBD}
\providecommand{\divthirteenC}{\TBD}
\providecommand{\divthirteenFirstErrTwo}{\TBD}
\providecommand{\divthirteenNitems}{\TBD}
\providecommand{\divthirteenPredYes}{\TBD}
\providecommand{\divthirteenTraceOK}{\TBD}
\providecommand{\divthirteenWrongN}{\TBD}
\providecommand{\divthirteenWrongZero}{\TBD}
\providecommand{\divthreeBAci}{\TBD}
\providecommand{\divthreeBAdiff}{\TBD}
\providecommand{\divthreeBase}{\TBD}
\providecommand{\divthreeC}{\TBD}
\providecommand{\mDiveleven}{\TBD}
\providecommand{\mDivseven}{\TBD}
\providecommand{\mDivsevenSix}{\TBD}
\providecommand{\mDivthirteen}{\TBD}
\providecommand{\mDivthirteenMid}{\TBD}
\providecommand{\mDivtwo}{\TBD}
\providecommand{\pDiveleven}{\TBD}
\providecommand{\pDivseven}{\TBD}
\providecommand{\pDivsevenSix}{\TBD}
\providecommand{\pDivthirteen}{\TBD}
\providecommand{\preregCommit}{\TBD}
\providecommand{\preregDate}{\TBD}
\providecommand{\preregTime}{\TBD}
\providecommand{\primeAsd}{\TBD}
\providecommand{\primeBAci}{\TBD}
\providecommand{\primeBAdiff}{\TBD}
\providecommand{\primeBAplarge}{\TBD}
\providecommand{\primeBsd}{\TBD}
\providecommand{\primeCompThreeSevenA}{\TBD}
\providecommand{\primeCompThreeSevenN}{\TBD}
\providecommand{\primeCsOneComps}{\TBD}
\providecommand{\primeCsOnePrimes}{\TBD}
\providecommand{\primeCsd}{\TBD}
\providecommand{\sEightNineTime}{\TBD}
\providecommand{\sFiveClarTime}{\TBD}
\providecommand{\sSevenTime}{\TBD}
\providecommand{\stepLocalAcc}{\TBD}
\providecommand{\stepLocalPowFour}{\TBD}
\providecommand{\validAsd}{\TBD}
\providecommand{\validBAci}{\TBD}
\providecommand{\validBAdiff}{\TBD}
\providecommand{\validBsd}{\TBD}
\providecommand{\validCsd}{\TBD}
\providecommand{\verdictSEight}{\TBD}
\providecommand{\verdictSNine}{\TBD}
\providecommand{\verdictSSeven}{\TBD}
\providecommand{\verdictSSevenD}{\TBD}

% added after an internal review pass run with AI agents: fallbacks for macros generated by make_numbers.py
\providecommand{\SsevenHeld}{\TBD}
\providecommand{\divelevenBAdOne}{\TBD}
\providecommand{\divelevenBAdZero}{\TBD}
\providecommand{\divsevenAnFiveForty}{\TBD}
\providecommand{\divsevenAnOneEighty}{\TBD}
\providecommand{\divsevenAnSixty}{\TBD}
\providecommand{\divsevenBnNinety}{\TBD}
\providecommand{\divsevenBnOneEighty}{\TBD}
\providecommand{\divsevenBnSixty}{\TBD}
\providecommand{\divsevenBnTwoSeventy}{\TBD}
\providecommand{\divsevenBpsOne}{\TBD}
\providecommand{\divsevenBpsZero}{\TBD}
\providecommand{\divsevenDsOne}{\TBD}
\providecommand{\divsevenDsZero}{\TBD}
\providecommand{\divsevenSixBsOne}{\TBD}
\providecommand{\divsevenSixBsZero}{\TBD}
\providecommand{\divsevenSixBsd}{\TBD}
\providecommand{\divthirteenAnOneEighty}{\TBD}
\providecommand{\divthirteenAnThreeSixty}{\TBD}
\providecommand{\divthirteenBAdOne}{\TBD}
\providecommand{\divthirteenBAdZero}{\TBD}
\providecommand{\divthirteenBnFiveForty}{\TBD}
\providecommand{\divthirteenBnOneEighty}{\TBD}
\providecommand{\divthirteenBnThreeSixty}{\TBD}
\providecommand{\divthreeBAdOne}{\TBD}
\providecommand{\divthreeBAdZero}{\TBD}
\providecommand{\halfBdivsevenBase}{\TBD}
\providecommand{\halfBprimeA}{\TBD}
\providecommand{\halfBprimeB}{\TBD}
\providecommand{\halfBprimeBase}{\TBD}
\providecommand{\mDivsevennNinety}{\TBD}
\providecommand{\mDivsevennTwoSeventy}{\TBD}
\providecommand{\mDivthirteennFiveForty}{\TBD}
\providecommand{\pDivsevenBp}{\TBD}
\providecommand{\pDivsevennNinety}{\TBD}
\providecommand{\pDivsevennOneEighty}{\TBD}
\providecommand{\pDivsevennTwoSeventy}{\TBD}
\providecommand{\pDivthirteennFiveForty}{\TBD}
\providecommand{\pDivthirteennOneEighty}{\TBD}
\providecommand{\pDivthirteennThreeSixty}{\TBD}
\providecommand{\pDivthree}{\TBD}
\providecommand{\pDivtwo}{\TBD}
\providecommand{\primeBApsTwo}{\TBD}
\providecommand{\primeHardAagree}{\TBD}
\providecommand{\sevenBdivsevenBase}{\TBD}
\providecommand{\threeBbasedivsevenN}{\TBD}
\providecommand{\threeBdivsevenABasepsOne}{\TBD}
\providecommand{\threeBdivsevenABasepsZero}{\TBD}
\providecommand{\threeBdivsevenAsOne}{\TBD}
\providecommand{\threeBdivsevenAsZero}{\TBD}
\providecommand{\threeBdivsevenBBasepsOne}{\TBD}
\providecommand{\threeBdivsevenBBasepsZero}{\TBD}
\providecommand{\threeBdivsevenBase}{\TBD}
\providecommand{\threeBdivsevenBsOne}{\TBD}
\providecommand{\threeBdivsevenBsZero}{\TBD}
\providecommand{\threeBdivthirteenA}{\TBD}
\providecommand{\threeBdivthirteenB}{\TBD}
\providecommand{\threeBdivthirteenBase}{\TBD}
\providecommand{\threeBpDivseven}{\TBD}
\providecommand{\threeBprimeAsOne}{\TBD}
\providecommand{\threeBprimeAsZero}{\TBD}
\providecommand{\threeBprimeBase}{\TBD}
\providecommand{\threeBprimeBsOne}{\TBD}
\providecommand{\threeBprimeBsZero}{\TBD}
\providecommand{\threeBprimeHardA}{\TBD}
\providecommand{\threeBprimeHardAagree}{\TBD}
\providecommand{\threeBvalidBase}{\TBD}
\providecommand{\validBApsTwo}{\TBD}
\providecommand{\verdictSTen}{\TBD}

% added after an internal review pass run with AI agents: fallbacks for macros generated by make_numbers.py
\providecommand{\primeBApList}{\TBD}
\providecommand{\validBApList}{\TBD}

% added after an internal review pass run with AI agents: fallbacks for macros generated by make_numbers.py
\providecommand{\sevenBdivthirteenA}{\TBD}
\providecommand{\sevenBdivthirteenB}{\TBD}
\providecommand{\sevenBdivthirteenBase}{\TBD}
\providecommand{\sevenBprimeA}{\TBD}
\providecommand{\sevenBprimeB}{\TBD}
\providecommand{\sevenBprimeBase}{\TBD}
\providecommand{\sevenBvalidA}{\TBD}
\providecommand{\sevenBvalidB}{\TBD}
\providecommand{\sevenBvalidBase}{\TBD}
\providecommand{\sevenBvalidBerrN}{\TBD}
\providecommand{\sevenBvalidBerrTemplate}{\TBD}
\providecommand{\sevenBvalidBerrYes}{\TBD}
\providecommand{\threeBdivelevenA}{\TBD}
\providecommand{\threeBdivelevenB}{\TBD}
\providecommand{\threeBdivelevenBase}{\TBD}
\providecommand{\threeBdivsevenBp}{\TBD}
\providecommand{\threeBdivsevenD}{\TBD}
\providecommand{\threeBdivsevenS}{\TBD}
\providecommand{\threeBdivsevenSixA}{\TBD}
\providecommand{\threeBdivsevenSixB}{\TBD}
\providecommand{\threeBdivsevenSixBase}{\TBD}
\providecommand{\threeBdivthirteenAnThreeSixty}{\TBD}
\providecommand{\threeBdivthirteenBnThreeSixty}{\TBD}
\providecommand{\threeBdivthreeA}{\TBD}
\providecommand{\threeBdivthreeB}{\TBD}
\providecommand{\threeBdivthreeBase}{\TBD}
\providecommand{\threeBvalidS}{\TBD}
% S12 answer-only sweep and CoT-prompt evals (filled by make_numbers.py when the cells land)
\providecommand{\sweepAlrTwoEpTenSZero}{--}\providecommand{\sweepAlrTwoEpTenSOne}{--}
\providecommand{\sweepAlrTwoEpThirtySZero}{--}\providecommand{\sweepAlrTwoEpThirtySOne}{--}
\providecommand{\sweepAlrFiveEpTenSZero}{--}\providecommand{\sweepAlrFiveEpTenSOne}{--}
\providecommand{\sweepAlrFiveEpThirtySZero}{--}\providecommand{\sweepAlrFiveEpThirtySOne}{--}
\providecommand{\sweepAlrTwoEpTenSZeroYes}{--}\providecommand{\sweepAlrTwoEpTenSOneYes}{--}
\providecommand{\sweepAlrTwoEpThirtySZeroYes}{--}\providecommand{\sweepAlrTwoEpThirtySOneYes}{--}
\providecommand{\sweepAlrFiveEpTenSZeroYes}{--}\providecommand{\sweepAlrFiveEpTenSOneYes}{--}
\providecommand{\sweepAlrFiveEpThirtySZeroYes}{--}\providecommand{\sweepAlrFiveEpThirtySOneYes}{--}
\providecommand{\sweepAlrTwoEpTenTok}{--}\providecommand{\sweepAlrTwoEpThirtyTok}{--}
\providecommand{\sweepAlrFiveEpTenTok}{--}\providecommand{\sweepAlrFiveEpThirtyTok}{--}
\providecommand{\sweepAmax}{--}\providecommand{\sweepAmaxConfig}{--}\providecommand{\sweepAnonCollapse}{--}
\providecommand{\sweepDivthreeA}{--}\providecommand{\STwelveVerdict}{pending}\providecommand{\STwelveResult}{pending}
\providecommand{\divsevenAcotsZero}{--}\providecommand{\divsevenAcotsOne}{--}\providecommand{\divsevenAcot}{--}
\providecommand{\divthirteenBaseCot}{--}\providecommand{\divthirteenAcot}{--}\providecommand{\primeAcot}{--}
\providecommand{\divsevenAcotMaxTok}{--}\providecommand{\divsevenBaseCotMedTok}{--}\providecommand{\divsevenAcotYes}{--}
\providecommand{\sevenBvalidBsOne}{--}\providecommand{\threeBdivsevenBnNinetysOne}{--}\providecommand{\threeBdivsevenBnTwoSeventysOne}{--}
\providecommand{\divsevenBnNinetysOne}{--}\providecommand{\divsevenBnTwoSeventysOne}{--}\providecommand{\divthirteenBnThreeSixtysOne}{--}\providecommand{\divthirteenBnFiveFortysOne}{--}
\providecommand{\divsevenBnNinetysZero}{--}\providecommand{\divsevenBnTwoSeventysZero}{--}\providecommand{\divthirteenBnThreeSixtysZero}{--}\providecommand{\divthirteenBnFiveFortysZero}{--}
\providecommand{\secondFamilyModel}{--}\providecommand{\famDivsevenAsZero}{--}\providecommand{\famDivsevenAsOne}{--}\providecommand{\famDivsevenBsZero}{--}\providecommand{\famDivsevenBsOne}{--}
\providecommand{\famDivsevenBAdiffsZero}{--}\providecommand{\famDivsevenBAdiffsOne}{--}\providecommand{\famDivsevenDsZero}{--}\providecommand{\famDivsevenBasesZero}{--}
\providecommand{\famPrimeAsZero}{--}\providecommand{\famPrimeBsZero}{--}\providecommand{\famValidAsZero}{--}\providecommand{\famValidBsZero}{--}\providecommand{\SElevenVerdict}{pending}\providecommand{\SElevenResult}{pending}
\providecommand{\threeBdivsevenAcot}{--}\providecommand{\threeBdivthirteenAcot}{--}\providecommand{\sevenBdivsevenAcot}{--}\providecommand{\sevenBdivthirteenAcot}{--}\providecommand{\bigAcotShortMin}{--}\providecommand{\sevenBdivsevenSsOne}{--}


\theoremstyle{definition}
\newtheorem{remark}{Remark}
\newcommand{\pp}{\,pp}
\newcommand{\nr}{--}   % cell not run (see table notes)
\def\NRToken{--}\def\PendingToken{pending}   % \ifx tests for macros that have not been generated yet
% figure if present (figures/ or paper root), otherwise a framed placeholder of the given height
\newcommand{\figorbox}[3]{%
  \IfFileExists{figures/#1}{\includegraphics[width=#2]{figures/#1}}{%
  \IfFileExists{#1}{\includegraphics[width=#2]{#1}}{%
  \fbox{\parbox[c][#3][c]{#2}{\centering\small\textit{Figure missing}}}}}}

\title{Perfect Labels for Free?\\ What Exact Answers and Exact Traces\\ Teach a Small Language Model}

% Submission is double blind; the author block is filled in for the camera-ready only.
\author{Anonymous authors\\ Paper under double-blind review}

%\iclrfinalcopy % camera-ready only
\begin{document}
\maketitle

\begin{abstract}
Exact verifiers (primality tests, SAT solvers, unit tests) label data without noise and supply the rewards of verifiable-reward post-training. Because every label is exact, a difference between two training regimes on the same inputs is a difference in the supervision format. We fine-tune Qwen2.5 models (0.5B to 7B; main size 1.5B, three seeds) with LoRA on exactly labeled decision tasks and compare supervision by the answer alone with a deterministic procedure trace followed by the answer. \emph{Exact answers teach shortcuts}: on 4-digit primality the answers follow ``the last digit is 1, 3, 7 or 9'' on \primeAagreeLastDigitLo--\primeAagreeLastDigitHi\% of test items, and on divisibility by 7 or 13 answer-only training stays at chance or constant at every size and configuration we ran ($n\le540$, including a logged sweep of four configurations, eight runs, none collapsing), while the same recipe learns divisibility by 3. \emph{Exact answers remove computation from the answers}: at 3B and 7B the zero-shot model already divides by 7; answer-only fine-tuning drops it to chance in both seeds, and under a chain-of-thought prompt the answer-only model still answers at once at every size tested (at least \bigAcotShortMin\% of items in three tokens or fewer) and stays at chance, while fine-tuning on its own verifier-filtered traces keeps the division (two seeds at each size); whether the capability survives is untested. \emph{Exact traces teach steps within a task}: accuracy rises with the number of traces and with how often each table entry was supervised, a count-based prediction across tasks failed because steps differ in difficulty, and only the input's own trace, written before the answer, helps. A second model family replicates every direction but misses one preset threshold. Predictions were logged in timestamped commits before their cells ran; four of twelve failed, among them every cross-task, cross-length and cross-family prediction, and all are reported.
\end{abstract}

\section{Introduction}
Verifiable-reward post-training rests on programs that decide correctness exactly: answer checkers for mathematics, unit tests for code, solvers for logic \citep{lambert2024tulu,deepseek2025r1}. Such verifiers are often described as a source of perfect labels for free, and a verified label is never wrong. Whether a model trained on those labels learns the computation the verifier performs, or a cheaper function that agrees with it on the training distribution, is a separate question, usually confounded with label noise, task difficulty and supervision format.

With exact labels the first confound is gone. We choose small decision tasks whose structure is known in advance (which surface rules approximate the label, which steps the natural procedure takes), so that we can tell which function the model learned, and what exact supervision of a given format teaches.

\emph{Answer-only} supervision (arm A) teaches a shortcut (Figure~\ref{fig:story}b). On 4-digit primality A's answers coincide with ``the last digit is 1, 3, 7 or 9'' on \primeAagreeLastDigitLo--\primeAagreeLastDigitHi\% of test items, and where no surface rule exists (divisibility by 7 or 13) A stays at chance at every training size we ran. At 3B and 7B, where the zero-shot model already divides correctly, answer-only fine-tuning removes that computation from the answers in both seeds, while fine-tuning on the model's own verifier-filtered traces keeps it (Figure~\ref{fig:story}a); whether the division is lost or only suppressed is left open. \emph{Trace} supervision (arm B: a deterministic procedure trace, then the answer) teaches the procedure once its steps are supervised often enough (Figure~\ref{fig:map}a): on divisibility by 7, accuracy rises from \divsevenBnSixty\% at $n=\nTrainSmall$ training examples to \divsevenBnTwoSeventy\% at 270 (means over three seeds), and only the input's own trace, placed before the answer, does this (Figure~\ref{fig:story}c). It is not a free improvement: at 1.5B the primality trace hurts, on argument validity answers do as well, and at 3B and $n=180$ the long-division trace underperforms the zero-shot model.

\begin{figure}[t]
\centering
\figorbox{fig_story.pdf}{\linewidth}{5.2cm}
\caption{(a) Divisibility by 7: zero-shot accuracy, answer-only fine-tuning (A), trace fine-tuning (B) and fine-tuning on the model's own verifier-filtered traces (S) at 1.5B, 3B and 7B; points are seeds, bars means. Where the base model already divides (3B \threeBdivsevenBase\%, 7B \sevenBdivsevenBase\%), A drops it to chance (3B \threeBdivsevenAsZero/\threeBdivsevenAsOne, 7B \sevenBdivsevenAsZero/\sevenBdivsevenAsOne) and S keeps it (3B \threeBdivsevenSsZero/\threeBdivsevenSsOne, 7B \sevenBdivsevenSsZero/\sevenBdivsevenSsOne). (b) Primality: accuracy of A and B on the \storyPrimeRuleRightN{} test items where the rule ``last digit is 1, 3, 7 or 9'' is right (circles) and on the \storyPrimeRuleWrongN{} composites where it is wrong (crosses); at 1.5B and 3B, A scores \storyPrimeAconsLo--\storyPrimeAconsHi\% on the first group and \storyPrimeAinconsLo--\storyPrimeAinconsHi\% on the second; 7B seed 0 is the collapsed seed that answers ``No''. (c) Controls on div7 at 1.5B: scrambled traces (C), a correct trace of a different number (D) and the trace placed after the answer (B$'$) stay at chance, and a zero-shot CoT prompt (\divsevenCot\%) or \nShots{} trace demonstrations (\divsevenFew\%) do not approach the model's own verifier-filtered traces (S) or the procedure trace (B). (d) The same arms against supervised completion tokens per epoch (whitespace tokens, as in the run records): under the shared configuration A never receives B's token budget, so B$-$A confounds format with tokens (the S12 sweep in \S\ref{sec:shortcut} closes part of that gap); C and D in (c) match B's tokens and stay at chance.}
\label{fig:story}
\end{figure}

\paragraph{Contributions.}
\begin{itemize}[leftmargin=*,itemsep=1pt,topsep=2pt]
\item Answer-only fine-tuning on exact labels learns the surface rule on primality (at 1.5B and 3B and in one of two 7B seeds), stays at chance where no surface rule fits (div7, div13) at every size and configuration we ran ($n\le540$; a logged sweep of four configurations, eight runs, none collapsing, S12) and, at 3B and 7B, replaces the base model's working division with a guess in both seeds; under a CoT prompt the answer-only model still answers at once and stays at chance at 1.5B, 3B and 7B (whether the capability survives is untested), while the model's own verifier-filtered traces keep the division (\S\ref{sec:shortcut}, Table~\ref{tab:scale}).
\item Trace supervision teaches the procedure within a task once its steps are supervised often enough: accuracy rises with training traces and with per-entry coverage, while a count-based prediction across tasks failed, exposing step difficulty as a confound of supervision frequency (\S\ref{sec:traces}).
\item Controls and exact audits: the trace must be the input's own computation, placed before the answer, so trace semantics are load-bearing in fine-tuning where \citet{valmeekam2026beyond} found them dispensable when training from scratch (\S\ref{sec:controls}); trace-trained models also write shortcuts behind unfaithful steps, which the audits catch (\S\ref{sec:traces}, \S\ref{sec:scale}); every failed prediction is reported (Table~\ref{tab:prereg}).
\end{itemize}

\section{Related work}
\paragraph{Intermediate steps as supervision.} Scratchpads \citep{nye2021show} and chain-of-thought (CoT) prompting \citep{wei2022chain} improve multi-step computation, mainly on mathematical and symbolic tasks \citep{sprague2025cot}. Data format changes the sample efficiency of arithmetic learning \citep{lee2024teaching}, fine-tuned models can learn arithmetic with decomposed steps \citep{liu2023goat}, and GCD predictions follow rules built from divisors of the number base \citep{charton2024gcd}, the kind of surface rule we probe for; procedure cloning \citep{yang2022procedure} imitates the intermediate computation as our arm B does, and surface form limits arithmetic in pretrained transformers \citep{nogueira2021investigating}. \citet{lee2024teaching} also fine-tune pretrained models and find that detailed scratchpads beat plain answers; their labels are exact too. What is new here is not the labels but the instruments: a named surface rule and a fitted probe that say \emph{what} answer-only training learned, controls that give a correct trace of the wrong number or place the trace after the answer, and the observation that answer-only fine-tuning removes a computation the base model already performs. Errors compound across steps \citep{dziri2023faith}; filler tokens can replace a chain of thought on some tasks \citep{pfau2024dot}, and our scrambled control keeps the tokens but removes their order. Reasoning helps when training data are local, so that intermediate variables are seen together often \citep{prystawski2023why}; our dose results are a within-task instance of that locality account, and they show where it stops predicting across tasks. Distilling rationales into small models \citep{hsieh2023distilling,ho2023teachers,magister2023teaching} trains on traces before answers; placing explanations after the answer \citep{lampinen2022explanations} and internalizing steps \citep{deng2023implicit,deng2024explicit} are the alternatives our answer-first arm B$'$ probes, and on div7 it does not transfer the procedure. Transformers can also learn multi-hop composition implicitly after long training \citep{wang2024grokked}; our answer arm, at a few hundred examples, does not.

\paragraph{Theory.} Parity-like targets give gradient-based learners almost no signal from end labels \citep{shalevshwartz2017failures}. Sub-task decomposition makes otherwise unlearnable tasks learnable \citep{wies2023subtask}; CoT raises the expressivity of bounded-depth transformers \citep{feng2023towards,li2024chain,merrill2024expressive}, next-token predictors trained on CoT data are universal learners \citep{malach2024autoregressive}, and CoT induces sparse, sample-efficient dependence \citep{wen2025sparse}. Transformers can represent automata through shallow shortcuts \citep{liu2023shortcuts}. The globality degree of \citet{abbe2024far} bounds what transformers learn efficiently; teacher-forced intermediate parities make parity learnable \citep{kim2025parity}. Appendix~\ref{app:automaton} uses these results as intuition only.

\paragraph{Traces and verifiers in post-training.} Transformers trained from scratch on verifiable traces depend less on trace semantics than usually assumed \citep{valmeekam2026beyond}, and CoT explanations can be unfaithful to the computation behind the answer \citep{turpin2023unfaithful}; our exact audits measure this in fine-tuned traces. Process supervision rewards steps rather than outcomes \citep{lightman2024verify}; program-generated logic corpora improve reasoning \citep{morishita2024logic}. STaR \citep{zelikman2022star} and rejection-sampling fine-tuning \citep{yuan2023scaling} train on self-generated rationales that reach verified answers, as our arm S does; RL with verifiable rewards drives recent reasoning models \citep{lambert2024tulu,deepseek2025r1}. Fine-tuning on model-generated positives can amplify spurious correlations \citep{setlur2024rl}; SFT and RL generalize differently \citep{chu2025sft}; RLVR mostly reweights what the base model already samples \citep{yue2025rl}. Our answer arm is an instance of shortcut learning \citep{geirhos2020shortcut}; exact labels do not prevent it.

\section{Setup}
\label{sec:setup}
\paragraph{Tasks.} Each task has an exact decision procedure, so labels are correct by construction; arms see the same prompts and labels and differ only in the completion (Tables~\ref{tab:tasks} and~\ref{tab:arms}). \textbf{prime}: is a 4-digit $n$ prime? The trace performs trial division by the primes up to $\lfloor\sqrt n\rfloor$, one ``$n \bmod p = r$.'' step per prime, stopping at the first divisor. \textbf{div$d$}: is $n$ divisible by $d$? The trace performs long division, $r\leftarrow(10r+x_t)\bmod d$ for each digit $x_t$, for $d\in\{2,3,7,11,13\}$ with 4 digits and $d=7$ with 6. \textbf{valid}: does a propositional conclusion follow from its premises (Z3 \citep{demoura2008z3})? Arguments instantiate \nForms{} forms shared between training and test, so the task is surface-learnable by design (Appendix~\ref{app:logic}). Test sets have \nTest{} items (\nTestClass{} per class), are problem-disjoint from training and identical across arms; training sets are balanced samples of size $n$ drawn with the run's seed.

\begin{table}[t]
\caption{Tasks. $k$: steps in the trace (prime: for primes; a composite's trace stops at its smallest factor). Tokens: mean whitespace tokens per trace. $m$: supervised transitions per (remainder, digit) table entry at $n=\nTrain$.}
\label{tab:tasks}
\centering\small
\begin{tabular}{llccc}
\toprule
task & trace procedure & $k$ & tokens & $m$ \\
\midrule
prime & trial division by the primes $\le\lfloor\sqrt n\rfloor$ & \primeStepsMin--\primeStepsMax & \primeTraceTok & -- \\
div$d$, 4 digits & long division, $r\leftarrow(10r+x_t)\bmod d$ & \divsevenSteps & \divsevenTraceTok & \mDivtwo{} ($d$=2) \dots\ \mDivthirteen{} ($d$=13) \\
div7, 6 digits & same, $d=7$ & \divsevenSixSteps & \divsevenSixTraceTok & \mDivsevenSix \\
valid & Z3 countermodel or unsatisfiable core & certificate & \validTraceTok & -- \\
\bottomrule
\end{tabular}
\end{table}

\begin{table}[t]
\caption{Arms (completions trained on; the prompt is shared) and the labels of the tests. H1--H4: endpoints of a first preregistration (prior grid). S1--S12: timestamped predictions of the second file (Table~\ref{tab:prereg}).}
\label{tab:arms}
\centering\footnotesize\setlength{\tabcolsep}{4pt}
\begin{tabular}{@{}lp{0.86\linewidth}@{}}
\toprule
arm & completion \\
\midrule
base & none (zero-shot); also a zero-shot CoT prompt and \nShots{} trace demonstrations \\
A & ``Answer: Yes'' or ``Answer: No'' \\
B & the procedure trace, a newline, the answer \\
C & B's traces permuted across examples and token-shuffled; the answer keeps the true label \\
D (S8) & a fluent, correct trace of a \emph{different} number with the same label \\
B$'$ & the answer first, then the trace \\
S (S10) & the base model's own sampled solution, kept when the exact gate accepts its final answer (STaR-style) \\
A-full & answers on the whole training pool \\
\bottomrule
\end{tabular}
\end{table}

\paragraph{Training and evaluation.} LoRA \citep{hu2022lora} with $r=\loraRank$, $\alpha=\loraAlpha$, dropout \loraDropout{} on the attention projections of Qwen2.5-Instruct \citep{qwen2024qwen25} (0.5B, 1.5B, 3B, 7B); AdamW, learning rate $\learningRate$, batch size 1, \nEpochs{} epochs, bf16; loss on completion tokens only; one configuration for every arm, not tuned per arm. Evaluation is greedy generation with one answer extractor for every arm; an unparsed or truncated answer counts as wrong. Details and generation caps: Appendix~\ref{app:setup}.

\paragraph{Statistics and timestamped predictions.}
\label{sec:prereg}
We lead with the difference of two arms and its 95\% bootstrap interval, resampling items and seeds jointly, and give the number of seeds; with one to three seeds per cell these intervals are item-level, and the paper makes no seed-level inference; per-seed exact McNemar tests \citep{mcnemar1947note}, Holm-corrected \citep{holm1979simple} across tasks within a comparison, model and seed at $\alpha=\alphaLevel$, are secondary, because they count test items, not training runs. A first preregistration, written before a prior grid of \nPriorRuns{} runs according to the history of the repository it was written in, fixed H1--H4. A second file, committed at \preregTime{} on \preregDate{} before any of its runs, fixes S1--S6; its append-only decision log adds S7--S10, each committed before any cell it governs had started (Appendix~\ref{app:prereg}). In short: H1--H4 were fixed before the prior grid; S1--S6 before any run of this study; S7--S10 were logged during the study, each before its own cells, and we call them \emph{timestamped} rather than preregistered; S3 and S6 were effectively settled before registration. S1--S10 were committed between \preregTime{} and 21:46:45 on \preregDate{}, with other cells running concurrently; S11 and S12 at 01:35:27 and 02:41:32 on 2026-09-26, each before any of its cells ran. S7--S12 were proposed by the AI assistant after earlier results were in; S12 answered a question raised by an internal review pass run with AI agents, not by peer review. The tests at risk were S1, S2, S4, S5 and S7--S10. Of the ten, S1--S4, S6 and S8 passed, S5, S9 and S10 failed, and S7 met its 4-of-5 rule with three of four informative cells; two later entries, S11 (a second model family, \S\ref{sec:scale}) and S12 (an answer-only sweep, \S\ref{sec:shortcut}), were logged in the final hours; S11 \SElevenVerdict{} on its threshold and S12 \STwelveVerdict{} (Table~\ref{tab:prereg}, Appendix~\ref{app:tables}). The surface probe and every analysis labelled post hoc belong to neither file.

\section{Results}
\label{sec:results}

\begin{table}[t]
\caption{Accuracy (\%) at $n=\nTrain$, Qwen2.5-1.5B-Instruct, \nTest{} test items per task; mean over seeds $\pm$ SD across seeds (seeds per cell: Table~\ref{tab:runs}; no SD: one seed). probe: the best surface probe fitted on A's training items (\S\ref{sec:shortcut}). B$-$A: points, with a 95\% bootstrap interval over items and seeds. max $p$: largest per-seed exact McNemar $p$ (B vs.\ A); in parentheses where A $>$ B. \nr: not run.}
\label{tab:main}
\centering\footnotesize\setlength{\tabcolsep}{2.2pt}
\begin{tabular}{lcccccll}
\toprule
task & base & probe & A & B & C & B$-$A [95\% CI] & max $p$ \\
\midrule
prime & \primeBase & \primeProbe & \textbf{\primeA}\primeAsd & \primeB\primeBsd & \primeC\primeCsd & \primeBAdiff{} \primeBAci & \primeBApMaxCol \\
div7 & \divsevenBase & \divsevenProbe & \divsevenA\divsevenAsd & \textbf{\divsevenB}\divsevenBsd & \divsevenC\divsevenCsd & \divsevenBAdiff{} \divsevenBAci & \divsevenBApMaxCol \\
valid & \validBase & \validProbe & \textbf{\validA}\validAsd & \validB\validBsd & \validC\validCsd & \validBAdiff{} \validBAci & \validBApMaxCol \\
div3 & \divthreeBase & \divthreeProbe & \divthreeA\divthreeAsd & \divthreeB\divthreeBsd & \nr & \divthreeBAdiff{} \divthreeBAci & \divthreeBApMaxCol \\
div11 & \divelevenBase & \divelevenProbe & \divelevenA\divelevenAsd & \divelevenB\divelevenBsd & \nr & \divelevenBAdiff{} \divelevenBAci & \divelevenBApMaxCol \\
div13 & \divthirteenBase & \divthirteenProbe & \divthirteenA\divthirteenAsd & \divthirteenB\divthirteenBsd & \divthirteenC\divthirteenCsd & \divthirteenBAdiff{} \divthirteenBAci & \divthirteenBApMaxCol \\
div7 (6 digits) & \divsevenSixBase & \divsevenSixProbe & \divsevenSixA\divsevenSixAsd & \divsevenSixB\divsevenSixBsd & \nr & \divsevenSixBAdiff{} \divsevenSixBAci & \divsevenSixBApMaxCol \\
\bottomrule
\end{tabular}
\end{table}

\subsection{Exact answer labels teach the surface and erase computation}
\label{sec:shortcut}
\paragraph{The primality shortcut (descriptive, post hoc).} A reaches \primeA\% on prime (\primeASeeds), and its answers coincide with ``the last digit is 1, 3, 7 or 9'' on \primeAagreeLastDigitLo--\primeAagreeLastDigitHi\% of test items; that rule alone scores \ruleLastDigit\%. A classifies \primeAprimesCorrect\% of primes and \primeAcompCorrect\% of composites correctly, and of the \primeCompLastDigitN{} test composites ending in 1, 3, 7 or 9 it gets \primeCompLastDigitA{} right in \primeASeeds, including \primeCompThreeSevenN{} with a factor 3 or 7 that any sieve would catch. On a dedicated set of \primeHardNneg{} odd composites without a factor $\le 7$ and as many primes (S6, \verdictSSix, settled in advance; the file's rule, ``odd and no factor $\le7$'', coincides with the last-digit rule on that set), A reaches \primeHardA\% and agrees with the rule on \primeHardAagree\% of items. Split by the rule itself (Figure~\ref{fig:story}b), A is right on \storyPrimeAconsLo--\storyPrimeAconsHi\% of the \storyPrimeRuleRightN{} test items where the rule is right and on \storyPrimeAinconsLo--\storyPrimeAinconsHi\% of the \storyPrimeRuleWrongN{} composites where it is wrong (\storyPrimeAcells{} cells at 1.5B and 3B). We named the rule after seeing the prior grid, so this analysis describes rather than tests.

\paragraph{Answers track surface probes.} We fit logistic regression, a one-hidden-layer MLP and $k$-nearest neighbours on the same \nTrain{} training items as arm A, on \nProbeFeatures{} one-hot digit-by-position features alone or with engineered divisibility features (last digit even or 0/5; digit sum mod 3 and 9; alternating digit sum mod 11), or for arguments on atom-renamed word and character $n$-grams, and choose the probe by 5-fold cross-validation; on div2, div3 and div11 the engineered features contain the answer, so there the probe is one-hot only. The probe is a property of the data, a lower bound on what a surface learner can decode from the training items, not a property of the model, and it is informative only where it sits between chance and ceiling. A tracks the probe (Figure~\ref{fig:map}c): on prime it lies between the one-hot linear probe (\primeProbeOneHot\%) and the best probe (\primeProbe\%), on div2 both are at ceiling (\divtwoA\% and \divtwoProbe\%) and on every other divisibility task both are at chance, and on valid both are near ceiling, where agreement is uninformative.

\paragraph{More answers do not buy the program.} On prime, A reaches \primeAsmall\%, \primeA\% and \primeAlarge\% at $n=\nTrainSmall$, \nTrain{} and \nTrainLarge{}, and \primeAtok\% on the full pool. On div7 it stays at chance at every $n$ (\divsevenAnSixty\%, \divsevenAnOneEighty\%, \divsevenAnFiveForty\%) and on div13 likewise (\divthirteenAnOneEighty\% and \divthirteenAnThreeSixty\% at \nTrain{} and 360), and at $n=540$ (div7) and $n=360$ (div13) A collapses to answering ``No'' under the untuned configuration (Limitations); on div7 it falls below the zero-shot base (\divsevenBase\%; A$-$base \divsevenBootABaseDiff{} \divsevenBootABaseCi{} points over \divsevenBootABaseSeeds{} seeds; significant after Holm correction in \divsevenABaseHolmSig{} seeds, $p_{\mathrm{Holm}}=\divsevenABaseHolmsZero$, $\divsevenABaseHolmsOne$, $\divsevenABaseHolmsTwo$). With labels permuted across items, the earlier study's answer arm sat at \shufMin--\shufMax\% on all seven of its tasks (Appendix~\ref{app:tables}): correct labels drive A, and on primality what they buy is the rule.

\paragraph{Is A's failure on div7 the configuration? (S12, timestamped).} Every A cell above uses the one shared configuration, under which A sees \tokRatioLow--\tokRatioHigh{} times fewer supervised tokens than B (counted with the model's tokenizer) and sometimes collapses to a constant answer. A sweep logged before it ran (S12) trains A on div7 at $n=\nTrain$ with learning rates $2\times10^{-5}$ and $5\times10^{-5}$ and 10 or 30 epochs, two seeds each, which gives A \sweepAlrTwoEpTenTok{} to \sweepAlrFiveEpThirtyTok{} supervised whitespace tokens in total (the unit of the run records), so that at 30 epochs A sees as many whitespace tokens as B sees per epoch on div7; the rule was that no configuration reaches 65\%, at least one does not collapse, and div3 A reaches 90\%. \ifx\sweepAmax\NRToken The cells were running at submission; the result will be reported in the supplementary material and here.\else Result: \STwelveResult; S12 \STwelveVerdict.\fi

\paragraph{Answer-only fine-tuning removes the base model's division (3B, 7B).} The 3B and 7B zero-shot models already decide divisibility: \threeBdivsevenBase\% and \threeBdivthirteenBase\% on div7 and div13 at 3B, \sevenBdivsevenBase\% and \sevenBdivthirteenBase\% at 7B. The 3B base model writes a direct division line (``$2740 \div 7 = 391.43\ldots$'') on \threeBbasedivsevenDirect\% of its div7 answers. Answer-only fine-tuning removes it (Figure~\ref{fig:story}a): div7 A falls to \threeBdivsevenA\% at 3B (A$-$base \threeBdivsevenBootABaseDiff{} \threeBdivsevenBootABaseCi) and \sevenBdivsevenA\% at 7B (\sevenBdivsevenBootABaseDiff{} \sevenBdivsevenBootABaseCi), and div13 A to \threeBdivthirteenA\% and \sevenBdivthirteenA\% (two seeds each; Table~\ref{tab:runs}). The mechanism is not a learned rule: the 3B answer-only runs collapse to a constant answer and the 7B runs guess without structure, and no divisibility A run agrees with any surface rule beyond $\kappa=0.24$ (post hoc audit, Appendix~\ref{app:tables}). Item by item, A keeps only \erThreeBdivsevenAkeep\% (3B) and \erSevenBdivsevenAkeep\% (7B) of the div7 items the base solved, the share a guesser would keep, while S keeps \erThreeBdivsevenSkeep\% and \erSevenBdivsevenSkeep\% (7B: seed 0) (post hoc paired audit, Appendix~\ref{app:tables}).

\paragraph{The model's own verifier-filtered traces keep it (arm S).} Fine-tuning on the model's own verifier-filtered solutions (arm S) keeps the computation: at 3B, S reaches \threeBdivsevenS\% on div7 (\threeBdivsevenSsZero, \threeBdivsevenSsOne), its errors fall on the base's own failures, its few fixes of those failures are not significant after Holm correction, and its kept traces are right (the post hoc re-audit parses \threeBselfParsedExt\% and finds \threeBselfCorrectParsedExt\% fully correct; the training script's narrower audit parsed \threeBselfParsed\%). At 7B, S gives \sevenBdivsevenS\% (two seeds), level with the base (seed-0 discordant pairs \sevenBdivsevenSBaseDisc). The foreign trace's dependence on size and the 3B dose replication are in Appendix~\ref{app:tables}.

\paragraph{A chain-of-thought prompt does not bring it back.} \ifx\divsevenAcot\NRToken Whether the answer-only adapter has lost the division or only stopped emitting it is untested at 3B and 7B: we did not prompt those A-tuned models for a chain of thought, and removing the intermediate tokens removes serial computation by construction \citep{merrill2024expressive}.\else Removing the intermediate tokens removes serial computation by construction \citep{merrill2024expressive}, so we asked whether the answer-only model still divides when asked for a chain of thought. It does not, and it does not try: at 1.5B, every answer-only adapter we tested (div7, two seeds; div13; prime) answers the CoT prompt in at most \divsevenAcotMaxTok{} tokens on every item, where the base model writes a median of \divsevenBaseCotMedTok{} tokens under the same prompt. The div7 adapters score \divsevenAcotsZero\% and \divsevenAcotsOne\%, near-constant answers in opposite directions (base with the prompt \divsevenCot\%)\ifx\divthirteenAcot\NRToken\else; the div13 adapter stays constant at \divthirteenAcot\% (the base reasons under the prompt but is near chance on div13, \divthirteenBaseCot\%); the primality adapter keeps its shortcut accuracy (\primeAcot\%)\fi. At 3B and 7B (seed 0, one run per cell) the answer-only adapters do the same: at least \bigAcotShortMin\% of items answered in three tokens or fewer, div7 \threeBdivsevenAcot\% and \sevenBdivsevenAcot\%, div13 \threeBdivthirteenAcot\% and \sevenBdivthirteenAcot\%. Answer-only fine-tuning suppresses the reasoning behaviour itself; whether the capability underneath survives cannot be told from this test.\fi

\subsection{Exact traces teach steps seen often enough}
\label{sec:traces}
\begin{figure}[t]
\centering
\figorbox{fig_map.pdf}{\linewidth}{4.6cm}
\caption{(a) Accuracy against training examples $n$ for div7 and div13 at 1.5B: trace arm B (solid) and answer arm A (dashed; chance at every $n$); filled markers are means over two or three seeds, hollow markers single seeds. (b) Arm B per-step accuracy $q$ (probability that a step is right given that the previous one was; mean over seeds, bars show the range over seeds) against $m=kn/(10d)$, the supervised transitions per (remainder, digit) table entry, known before training; lines join the training sizes of div7 (orange) and div13 (violet); the 6-digit div7 cell (diamond) shares $m$ with div7 at $n=270$. Filled: cells predicted in advance by S7 or S9; hollow: cells observed first. Answer accuracy is in panel (a). (c) Answer arm A against the surface probe fitted on the same training items (one-hot features where engineered features encode the answer); dashed line: equality; points at chance or at ceiling are uninformative, and prime is the one informative point.}
\label{fig:map}
\end{figure}
\paragraph{Divisibility by 7.} At $n=\nTrain$, over three seeds, B reaches \divsevenB\% (\divsevenBsZero, \divsevenBsOne, \divsevenBsTwo) against A's \divsevenA\% (B$-$A \divsevenBAdiff{} \divsevenBAci), and C stays at \divsevenC\%. Seeds 1 and 2 are the preregistered replication (S1 \verdictSOne); a deterministic re-execution of seed 0 gave \divsevenBsZeroRerun\%. The gain holds on both classes (\divsevenByes\% on multiples, \divsevenBno\% on non-multiples).

\paragraph{Dose.} With the trace format fixed, accuracy rises with the number of training traces while A stays flat (Figure~\ref{fig:map}a, where hollow markers are single seeds). Div7 B goes from \divsevenBnSixty\% at $n=\nTrainSmall$ to \divsevenBnNinety\% at 90 (seed-0 per-step accuracy $q=\pDivsevennNinetysZero$), \divsevenBnOneEighty\% at \nTrain{} and \divsevenBnTwoSeventy\% at 270 (seed-0 $q=\pDivsevennTwoSeventysZero$), three seeds at each $n$ (per seed: Table~\ref{tab:runs}). Div13 B, with the same 4 steps, goes from \divthirteenBnOneEighty\% at \nTrain{} (seed-0 $q=\pDivthirteennOneEightysZero$) to \divthirteenBnThreeSixty\% at 360 and \divthirteenBnFiveForty\% at 540 (seed-0 $q=\pDivthirteennFiveFortysZero$), three seeds each.\footnote{Sensitivity check, not the preregistered measure: excluding the near-trivial first step ($r_1=x_1\bmod d$), seed-0 $q$ is \qAfterDivthirteenOneEighty{} instead of \qAllDivthirteenOneEighty{} for div13 at $n=180$ and \qAfterDivsevenNinety{} instead of \qAllDivsevenNinety{} for div7 at $n=90$.} At $n=\nTrain$, div13's failure is visible in its seed-0 generations: \divthirteenTraceOK\% of traces are correct and \divthirteenWrongZero{} of \divthirteenWrongN{} wrong traces end at $r=0$, so this seed answers ``yes'' \divthirteenPredYes\% of the time. In div7 generations (seed 0), digits and carried remainders are copied without error (\divsevenDigitCopy\%, \divsevenCarry\%); every error is a wrong lookup of $(10r+x)\bmod 7$.

\paragraph{Timestamped dose predictions (S7, S9): held within tasks, failed across tasks.} S7 predicted that $q$ follows $m=kn/(10d)$ in five cells and passes if four hold. The three within-task cells hold: (a) div13 at $n=540$ ($m=\mDivthirteennFiveForty$), $q=\pDivthirteennFiveFortysZero$ and accuracy \divthirteenBnFiveFortysZero\% (seeds 1 and 2, run later: \divthirteenBnFiveFortysOne\% and \divthirteenBnFiveFortysTwo\%); (b) div7 at $n=270$ ($m=\mDivsevennTwoSeventy$), $q=\pDivsevennTwoSeventysZero$; (c) div7 at $n=90$ ($m=\mDivsevennNinety$), $q=\pDivsevennNinetysZero$ and accuracy \divsevenBnNinetysZero\%. Cell (e), div3 and div2, \SsevenCellE{} only at ceiling ($q=\pDivthree$, \pDivtwo). (d) Div11 at $m=\mDiveleven$ reached $q=\pDiveleven$, above the predicted interval: $10\equiv-1\pmod{11}$ makes its step a subtraction. S9 (\verdictSNine): div13 at $n=360$ ($m=\mDivthirteenMid$) reached \divthirteenBnThreeSixtysZero\%, \divthirteenBnThreeSixtysOne\% and \divthirteenBnThreeSixtysTwo\% (seeds 0--2) against the required 80\%, while div7 at the similar $m=\mDivseven$ reached \divsevenBnOneEighty\%; the div13 step, $r'=(x-3r)\bmod 13$, is harder. S7 thus \verdictSSeven{} by its 4-of-5 rule (\SsevenHeld/5) only because cell (e) held at ceiling; of the four informative cells, the three within-task cells held and the cross-task cell failed, and that is how we report it. The supported reading is narrow: within a task more supervised transitions raise $q$; across tasks step difficulty shifts the curve, and $m$ does not transfer across tasks or sizes (post hoc fits of $q$ on $\ln m$: $R^2=\fitRsqOneFive$ at 1.5B, \fitRsqThreeB{} at 3B). S2 \verdictSTwo{} at 1.5B: B$-$A $\ge 15$ points in the two preregistered seeds for div3 (\divthreeBAdZero, \divthreeBAdOne) and div11 (\divelevenBAdZero, \divelevenBAdOne), not for div13 (\divthirteenBAdZero, \divthirteenBAdOne).

\paragraph{Coverage at the entry level (post hoc).} S7 counts supervised transitions per table entry on average; an audit written after the runs counts them per entry. For every step a trace-trained model wrote, let $c$ be the number of times that exact (remainder, digit) entry appeared in its own training traces. Over \covSteps{} steps from every arm-B cell with saved generations (all sizes, all divisibility tasks), lookup accuracy rises with $c$ from \covNTzero\% at $c=0$ to \covNTthirtythree\% at $c\ge33$ (Table~\ref{tab:coverage}, Appendix~\ref{app:tables}); within cells, the rank correlation between an entry's $c$ and its accuracy is positive in \covRhoPos{} of \covRhoCells{} cells (median $\rho=\covRhoMedian$, Wilcoxon $p=\covRhoP$; the draw of training items also sets their order and the initialization). Entries never seen in training are right in \covCzeroElevenK{} of \covCzeroElevenN{} div11 steps, consistent with a rule ($10\equiv-1\pmod{11}$), and in \covCzeroSevenK{} of \covCzeroSevenN{} div7 steps, consistent with lookup. At the same size, $n$ and $c$, div13 entries are learned more slowly than div7 entries (1.5B, $n=180$, $c=5$--8: \covDthirteenfiveeight\% against \covDsevenfiveeight\%; Table~\ref{tab:coverage}), which is the step-difficulty confound of S7 seen directly. Arm D shows the other half of what a trace teaches: its lookups are \covDlookup\% correct as table entries but only \covDlocal\% correct for the input's own digits, so D learned the table and not to read the number.


\paragraph{Length (S5): failed as preregistered.} The prediction for 6-digit div7 was B $\approx q^6=\divsevenSixPred\%$, with $q=\stepAcc$ measured on the seed-0 4-digit generations. The primary cell, B trained and tested on 6-digit traces, reached \divsevenSixBsZero\% and \divsevenSixBsOne\% in seeds 0 and 1, outside the 10-point tolerance: S5 \verdictSFive. Its own seed-0 $q$ is \pDivsevenSixsZero: at $k=6$ the same $n$ gives more transitions per entry ($m=\mDivsevenSix$). The secondary cell, the 4-digit model applied to 6-digit inputs, gives \divsevenTransferSixB\%, within the tolerance, but its own $q$ is \pDivsevenTransfer{} and only \divsevenTransferTraceOK\% of its traces are fully correct, so the match is not evidence for compounding (post hoc forms: Appendix~\ref{app:compound}).

\paragraph{Primality: the trace hurts at 1.5B.} B reaches \primeB\% against A's \primeA\% (\primeBSeeds; B$-$A \primeBootBADiff{} \primeBootBACi), and B vs.\ C is not significant ($p=\primeBCpsZero$, $\primeBCpsOne$). B classifies \primeBprimesCorrect\% of primes and \primeBcompCorrect\% of composites correctly. Class and trace length are confounded here: every prime needs the full sequence of trial divisions, and accuracy falls with the true number of divisions (\primeBdivBinOne, \primeBdivBinTwo, \primeBdivBinThree, \primeBdivBinFour{} and \primeBdivBinFive\% for 1, 2--3, 4--8, 9--16 and 17 or more). Scrambled C does not separate the two: seed 1 shows B's split, seed 0 the opposite (\primeCsZeroPrimes\% of primes, \primeCsZeroComps\% of composites). At $n=\nTrainLarge$, B reaches \primeBlarge\% against A's \primeAlarge\% (one seed). A trial-division step is a table indexed by $(n,p)$ with more than $10^5$ entries, so $m\ll1$ at these $n$. At 3B the trace arm scores higher (\threeBprimeB\%) without executing its trace: only \threeBprimeClaimAcc\% of its ``$n \bmod p = r$'' claims are correct, seed 0 agrees with ``no factor $\le 5$, so prime'' on \threeBprimeBnoFiveAgree\% of items, and on the hard set it scores \threeBprimeHardBsZero\% and \threeBprimeHardBsOne\%: a shortcut written in the form of a trace \citep[cf.][]{turpin2023unfaithful}.

\paragraph{Validity and cost.} On valid, answers match or beat traces (A \validA\%, B \validB\%, C \validC\%; B$-$A \validBootBADiff{} \validBootBACi; per seed \validBApList). B's completions contain \tokRatioLow--\tokRatioHigh{} times as many tokens as A's (Qwen tokenizer).

\subsection{Controls (div7, 1.5B)}
\label{sec:controls}
\paragraph{The trace must be the input's own computation.} D, a fluent and correct trace of a different number with the right label, gives \divsevenDsZero\% and \divsevenDsOne\% (S8 \verdictSEight): a correct trace of another number does not help. Scrambled C stays at \divsevenC\%. C controls for token count but not for interference: it also fails where A succeeds (valid \validC\% vs.\ \validA\%, prime \primeC\% vs.\ \primeA\%), so C alone does not show that the trace's information is what helps.

\paragraph{The trace must come before the answer.} B$'$, trained with the trace after the answer, gives \divsevenBpsZero\% and \divsevenBpsOne\%, like A, although it still writes a trace after answering with $q=\pDivsevenBp$.

\paragraph{Arm S at 1.5B (S10).} Sampling from the base model's CoT prompt and keeping completions the gate accepts covers \selfKept{} training items (keep rate \selfKeep\% of samples). S10 predicted div7 accuracy $\le 65\%$ in seeds 0 and 1, fewer than half of kept traces with all remainders correct, and valid $\ge 90\%$. Result: \STenResult, against \divsevenCot\% for the base model's CoT prompt; S10: \verdictSTen{} on its accuracy clauses. The failure is informative: arm S stayed above A in both seeds, and in seed 0 above the base model's own CoT prompt, so outcome-only filtering supplied traces that helped, against the prediction. The trace audit depends on its parser: the training script's audit counts \selfCorrectAll\% of kept traces as fully correct, while a post hoc re-audit that also reads other division formats counts \selfCorrectAllExt\%, so the trace clause of S10 holds under the preregistered parser and fails under the post hoc one; we score S10 by its accuracy clauses, which fail either way, and report the re-audit as a check rather than a verdict.


\section{Scale}
\label{sec:scale}
Table~\ref{tab:scale} and Figure~\ref{fig:arms} (Appendix~\ref{app:tables}) repeat the core cells at 0.5B, 3B and 7B. Three results replicate. Where A finds no surface rule (div7, div13, div11, 6-digit div7), answer-only training stays at chance at every size where it was run (Table~\ref{tab:scale}). The primality shortcut persists: at 3B, A reaches \threeBprimeA\% and agrees with the last-digit rule; at 7B one seed does the same (\sevenBprimeAsOne\%) while the other collapses to ``No'' on \sevenBprimeAsZeroNo{} of \nTest{} items, against B's \sevenBprimeB\% (\sevenBprimeBsZero, \sevenBprimeBsOne). Trace supervision beats answers on div7 at every size (S4 \verdictSFour; B$-$A \threeBdivsevenBootBADiff{} \threeBdivsevenBootBACi{} at 3B and \sevenBdivsevenBootBADiff{} \sevenBdivsevenBootBACi{} at 7B, two seeds each), and the div7 controls replicate: D gives \threeBdivsevenD\% at 3B and \sevenBdivsevenD\% at 7B, B$'$ \threeBdivsevenBp\% and \sevenBdivsevenBp\% (seeds and Yes-rates per cell in Table~\ref{tab:runs}; at 3B the D runs and one B$'$ run are constant-output); div11 at 7B behaves as at 1.5B (Table~\ref{tab:scale}).

Some cells move with size. At 3B, div3 becomes learnable from answers alone (A \threeBdivthreeA\%, against \divthreeA\% at 1.5B), so S2 fails at 3B, and the S9 threshold is met (B \threeBdivthirteenBnThreeSixtysZero\%, A \threeBdivthirteenAnThreeSixtysZero\%). At 7B the trace arm collapsed on valid in seed 0 (\sevenBvalidBsZero\%; seed 1 reached \sevenBvalidBsOne\%): all \sevenBvalidBerrYes{} of that seed's \sevenBvalidBerrN{} errors are invalid arguments answered ``yes'', \sevenBvalidBerrTemplate{} of them after a trace asserting a contradiction, the certificate's form without its check. The suppression result and its 3B dose replication are in \S\ref{sec:shortcut}.

\paragraph{A second model family (S11).} A prediction logged before its cells ran required, on \secondFamilyModel{} (the logged fallback; the first choice was gated), div7 B$-$A $\ge20$ points in two seeds and the grounding control D within 10 points of A. Result: \SElevenResult; S11 \SElevenVerdict. The direction replicates on every task (div7 A \famDivsevenAsZero/\famDivsevenAsOne, B \famDivsevenBsZero/\famDivsevenBsOne; prime A \famPrimeAsZero{} vs B \famPrimeBsZero; valid A \famValidAsZero{} vs B \famValidBsZero), with a smaller div7 gain at $n=\nTrain$ than in Qwen2.5-1.5B, as the dose results would predict; a larger-$n$ point in this family was not run.

An automaton view of the count $m$, and why it must fail across tasks, is in Appendix~\ref{app:automaton}.

\section{Discussion and limitations}
\label{sec:discussion}
\paragraph{What exact labels guarantee.} An exact verifier guarantees that no training label is wrong, not that the labels identify the verifier's function among those a model can fit. The rule A learned on primality calls every composite ending in 1, 3, 7 or 9 prime, including all \hardNegPop{} 4-digit composites without a factor $\le 7$. Shuffled labels cannot tell the verifier's function from a surface rule; inputs on which the rule is wrong can.

\paragraph{Answer-only fine-tuning suppresses the base model's test-time computation.} Labels are exact in every arm and we ran no noisy-label arm, so the effect comes from the supervision format. At 3B and 7B, answer supervision replaced a working computation with a guess, in both seeds; at 1.5B, A falls below the zero-shot base on div7 in the same direction (significant after Holm correction in \divsevenABaseHolmSig{} seeds). Fine-tuning on the model's own verifier-filtered solutions kept the computation (3B, \threeBdivsevenSSeedTxt; 7B, \sevenBdivsevenSSeedTxt). A pipeline that fine-tunes a capable model on verified answers alone should compare against the model's own zero-shot accuracy, not only against chance.

\paragraph{Supervision design: hypotheses for verifiable-reward pipelines.} Our results suggest, without establishing beyond these tasks: (1) fit a cheap surface probe on the verified data before training, and evaluate where it is wrong. (2) A correct trace is not enough: budget how often each step type is supervised, expect harder steps to need more, and check per-step accuracy after a small run. (3) The trace should be the input's own computation and precede the answer (div7 only). (4) Check the base model first: where its own steps are right, its verifier-filtered traces preserved them here (two seeds at 3B and at 7B). (5) Do not read a null from one training size: our \nTrainSmall-example null on div7 was the bottom of a dose curve. We ran no RL; arm S is the nearest experiment, and Appendix~\ref{app:rlvr} states the RLVR hypothesis these results suggest.

\paragraph{Limitations.} One model family for all but one test (the second family covers one size and one $n$), with 1.5B as the main size, and one hyperparameter configuration shared by all arms (LoRA $r=\loraRank$, learning rate $\learningRate$, \nEpochs{} epochs; Appendix~\ref{app:setup}), not tuned per arm; answer-only failure may reflect sample inefficiency under this configuration rather than impossibility, and A was tested only up to $n=\nTrainLarge$. B sees \tokRatioLow--\tokRatioHigh{} times more supervised tokens than A (tokenizer count; the run records count whitespace tokens), and the token-matched comparison (H3) was not run on div7. Binary tasks, \nTest{} test items, one to three seeds per cell, one trace format per task. Some base cells are lower bounds (capped generations) and some A seeds collapse to a constant answer (Appendix~\ref{app:setup}). Four of twelve timestamped predictions failed, $m$ is confounded with step difficulty, and the trace audits depend on the parser (\S\ref{sec:prereg}, \S\ref{sec:controls}).
\subsection*{AI use statement}
We used an AI assistant (Claude, Anthropic) for experiment code, running and logging experiments under our direction, statistics, drafting and editing, and reference checking; hypotheses S7--S12 were proposed by the assistant and committed before any of their cells ran (Appendix~\ref{app:prereg}). The authors reviewed all code, text and numbers and take full responsibility for this paper.

\subsection*{Reproducibility statement}
Code and data will be released: dataset generators with fixed seeds, the training and evaluation worker, queue runners, analysis scripts, both preregistration files with an anonymized timestamp manifest of their commits, and every run's result record with per-item correctness (local paths removed), from which each paired test can be recomputed; an anonymized copy is in the supplementary material. Section~\ref{sec:setup} and Appendices~\ref{app:setup} and~\ref{app:formats} give prompts, trace formats and hyperparameters. Every number in the text is a macro regenerated from the result files (\texttt{paper/exp/build\_all.py}).

\subsection*{Ethics statement}
The study uses synthetic arithmetic and logic data only and involves no human subjects or personal data.

\bibliography{refs}
\bibliographystyle{iclr2027_conference}

\appendix
\section{Full results}
\label{app:tables}
\paragraph{The foreign trace and the dose at 3B and 7B (not preregistered).} At 3B and 7B (not preregistered), the foreign long-division trace depends on size: at 3B it stays below the base on div7 (\threeBdivsevenB\%; B$-$base \threeBdivsevenBootBBaseDiff{} \threeBdivsevenBootBBaseCi) and div13 (\threeBdivthirteenB\%); at 7B it matches the base on div7 (\sevenBdivsevenB\%; \sevenBdivsevenBootBBaseDiff{} \sevenBdivsevenBootBBaseCi) but not on div13 (\sevenBdivthirteenB\%); at 1.5B, whose base scores \divsevenBase\%, it exceeds the base. More traces close the 3B gap, a within-task replication of the dose result above (not preregistered at 3B): div7 B rises from \threeBdivsevenBnNinetysZero\% at $n=90$ to \threeBdivsevenBnTwoSeventysZero\% at 270 in seed 0 (\threeBdivsevenBnNinetysOne\% and \threeBdivsevenBnTwoSeventysOne\% in seed 1), above the zero-shot model in one seed, and div13 B reaches \threeBdivthirteenBnFiveForty\% at $n=540$ while A stays at \threeBdivthirteenAnFiveForty\%.

\begin{table}[h]
\caption{Scale: zero-shot base / answers A / trace B at $n=\nTrain$, accuracy in \%, mean over seeds (seeds per cell: Table~\ref{tab:runs}). \nr: not run. $^\dagger$Lower bound: \threeBprimeBaseCap{} of \nTest{} generations hit the cap. Means that include a constant-output seed (7B prime A; 3B div13 A) are listed in Appendix~\ref{app:setup}.}
\label{tab:scale}
\centering\footnotesize\setlength{\tabcolsep}{4pt}
\begin{tabular}{lcccc}
\toprule
task & 0.5B & 1.5B & 3B & 7B \\
\midrule
div7 & \nr{} / \halfBdivsevenA{} / \halfBdivsevenB & \divsevenBase{} / \divsevenA{} / \divsevenB & \textbf{\threeBdivsevenBase} / \threeBdivsevenA{} / \threeBdivsevenB & \textbf{\sevenBdivsevenBase} / \sevenBdivsevenA{} / \sevenBdivsevenB \\
div13 & \nr{} / \halfBdivthirteenA{} / \halfBdivthirteenB & \divthirteenBase{} / \divthirteenA{} / \divthirteenB & \textbf{\threeBdivthirteenBase} / \threeBdivthirteenA{} / \threeBdivthirteenB & \textbf{\sevenBdivthirteenBase} / \sevenBdivthirteenA{} / \sevenBdivthirteenB \\
div3 & \nr{} / \halfBdivthreeA{} / \halfBdivthreeB & \divthreeBase{} / \divthreeA{} / \divthreeB & \nr{} / \threeBdivthreeA{} / \threeBdivthreeB & \nr \\
div11 & \nr & \divelevenBase{} / \divelevenA{} / \divelevenB & \nr{} / \threeBdivelevenA{} / \threeBdivelevenB & \nr{} / \sevenBdivelevenA{} / \sevenBdivelevenB \\
div7, 6 digits & \nr & \divsevenSixBase{} / \divsevenSixA{} / \divsevenSixB & \nr{} / \threeBdivsevenSixA{} / \threeBdivsevenSixB & \nr \\
prime & \nr{} / \halfBprimeA{} / \halfBprimeB & \primeBase{} / \primeA{} / \primeB & \threeBprimeBase$^\dagger$ / \threeBprimeA{} / \threeBprimeB & \sevenBprimeBase{} / \sevenBprimeA{} / \sevenBprimeB \\
valid & \nr{} / \halfBvalidA{} / \halfBvalidB & \validBase{} / \validA{} / \validB & \threeBvalidBase{} / \threeBvalidA{} / \threeBvalidB & \sevenBvalidBase{} / \sevenBvalidA{} / \sevenBvalidB \\
\bottomrule
\end{tabular}
\end{table}
\begin{table}[h]
\caption{Entry-level coverage (post hoc): lookup accuracy (\%) of trace-trained models by $c$, the number of times the exact (remainder, digit) entry appeared in the cell's training traces; non-trivial entries (remainder $\ne0$), pooled over every arm-B cell with saved generations, and div7 versus div13 at 1.5B and $n=180$ on the same bins.}
\label{tab:coverage}
\centering\footnotesize\setlength{\tabcolsep}{4pt}
\begin{tabular}{lcccccccc}
\toprule
$c$ & 0 & 1 & 2 & 3--4 & 5--8 & 9--16 & 17--32 & 33--64 \\
\midrule
all cells, non-trivial & \covNTzero & \covNTone & \covNTtwo & \covNTthreefour & \covNTfiveeight & \covNTninesixteen & \covNTseventeen & \covNTthirtythree \\
div7, 1.5B, $n=180$ & \covDsevenzero & \covDsevenone & \covDseventwo & \covDseventhreefour & \covDsevenfiveeight & \covDsevenninesixteen & \covDsevenseventeen & \covDseventhirtythree \\
div13, 1.5B, $n=180$ & \covDthirteenzero & \covDthirteenone & \covDthirteentwo & \covDthirteenthreefour & \covDthirteenfiveeight & \covDthirteenninesixteen & \covDthirteenseventeen & \covDthirteenthirtythree \\
\bottomrule
\end{tabular}
\end{table}

\begin{figure}[h]
\centering
\figorbox{fig_arms.pdf}{0.92\linewidth}{6cm}
\caption{Accuracy by arm and task at $n=\nTrain$; bars: 95\% bootstrap interval over items and seeds where a cell has two or more seeds, Wilson interval for a single seed; small dots are single seeds; one row per model with at least three tasks (the 0.5B cells are in Table~\ref{tab:scale}). Dashed line: majority class.}
\label{fig:arms}
\end{figure}

Table~\ref{tab:runs} lists every logged run of the arms in Table~\ref{tab:arms} and of B-ext, which uses reasoning traces sampled from a locally served 9B model and kept only when their final answer matched the verifier; only \primeBornN{} traces passed, so these runs are not comparable to B and are listed for completeness. Table~\ref{tab:shuffled} reports the earlier answer-only study with its shuffled-label control. Per-seed exact McNemar tests (B vs.\ A, 1.5B): prime $p=\primeBApsZero$, $\primeBApsOne$, $\primeBApsTwo$ (A $>$ B); div7 $\divsevenBApsZero$, $\divsevenBApsOne$, $\divsevenBApsTwo$; valid $\validBApsZero$, $\validBApsOne$, $\validBApsTwo$.

{\small\IfFileExists{tab_runs.tex}{% generated by paper/exp/make_numbers.py from the run files -- do not edit
\begin{longtable}{lllrrrrrr}
\caption{All logged runs of arms base, A, B, C, D, B$'$, S, A-full and B-ext (\nTest{} test items per cell). Cell: the test set, prefixed by the training task when they differ; [cot] and [fewshot4] mark prompting controls on the base model. SmolLM2: SmolLM2-1.7B-Instruct (S11). Train tokens: whitespace tokens in the completions. Cap: test generations that hit the generation cap (scored wrong; -- where not logged). Yes: share of test items answered Yes, from saved generations or recovered from per-item correctness; 0 or 100 marks a constant-output collapse.}\label{tab:runs}\\
\toprule
model & cell & arm & $n$ & seed & train tokens & cap & Yes (\%) & acc.\ (\%) \\
\midrule
\endfirsthead
\toprule
model & cell & arm & $n$ & seed & train tokens & cap & Yes (\%) & acc.\ (\%) \\
\midrule
\endhead
\bottomrule
\endfoot
0.5B & prime & A & 180 & 0 & 360 & 0 & 47 & 68.3 \\
0.5B & prime & B & 180 & 0 & 13{,}370 & 0 & 12 & 57.5 \\
\midrule
0.5B & div7 & A & 180 & 0 & 360 & 0 & 5 & 50.4 \\
0.5B & div7 & B & 180 & 0 & 10{,}800 & 0 & 34 & 79.2 \\
\midrule
1.5B & prime & base & 0 & 0 & 0 & -- & 28 & 59.6 \\
1.5B & prime & A & 60 & 0 & 120 & -- & 90 & 59.6 \\
1.5B & prime & A & 180 & 0 & 360 & -- & 67 & 82.1 \\
1.5B & prime & A & 180 & 1 & 360 & -- & 68 & 82.5 \\
1.5B & prime & A & 180 & 2 & 360 & 0 & 68 & 81.7 \\
1.5B & prime & A & 540 & 0 & 1{,}080 & -- & 75 & 75.0 \\
1.5B & prime & B & 60 & 0 & 4{,}355 & -- & 12 & 57.1 \\
1.5B & prime & B & 180 & 0 & 13{,}370 & -- & 11 & 57.9 \\
1.5B & prime & B & 180 & 1 & 13{,}260 & -- & 12 & 57.5 \\
1.5B & prime & B & 180 & 2 & 13{,}080 & 13 & 16 & 52.1 \\
1.5B & prime & B & 540 & 0 & 40{,}095 & -- & 28 & 68.8 \\
1.5B & prime & C & 180 & 0 & 12{,}561 & -- & 99 & 50.0 \\
1.5B & prime & C & 180 & 1 & 13{,}120 & -- & 12 & 54.2 \\
1.5B & prime & A-full & 1320 & 0 & 2{,}640 & -- & 69 & 80.8 \\
1.5B & prime & B-ext & 9 & 0 & 896 & -- & 20 & 62.1 \\
1.5B & prime & B-ext & 9 & 1 & 974 & -- & 12 & 60.0 \\
1.5B & prime & S & 180 & 0 & 18{,}790 & 1 & 41 & 67.1 \\
1.5B & prime$\to$prime\_hard & A & 180 & 0 & 0 & 0 & 99 & 49.6 \\
1.5B & prime[cot] & base & 0 & 0 & 0 & 8 & 28 & 62.9 \\
1.5B & prime[cot] & A & 180 & 0 & 0 & 0 & 67 & 82.5 \\
1.5B & prime[fewshot4] & base & 0 & 0 & 0 & 104 & 9 & 39.2 \\
\midrule
1.5B & div7 & base & 0 & 0 & 0 & -- & 72 & 62.5 \\
1.5B & div7 & A & 60 & 0 & 120 & -- & 88 & 50.0 \\
1.5B & div7 & A & 180 & 0 & 360 & -- & 24 & 52.1 \\
1.5B & div7 & A & 180 & 1 & 360 & -- & 82 & 50.8 \\
1.5B & div7 & A & 180 & 2 & 360 & -- & 94 & 47.9 \\
1.5B & div7 & A & 540 & 0 & 1{,}080 & -- & 0 & 50.0 \\
1.5B & div7 & B & 60 & 0 & 3{,}600 & -- & 73 & 50.0 \\
1.5B & div7 & B & 60 & 1 & 3{,}600 & 0 & 32 & 51.7 \\
1.5B & div7 & B & 60 & 2 & 3{,}600 & 0 & 48 & 53.3 \\
1.5B & div7 & B & 90 & 0 & 5{,}400 & 0 & 52 & 62.1 \\
1.5B & div7 & B & 90 & 1 & 5{,}400 & 0 & 21 & 56.7 \\
1.5B & div7 & B & 90 & 2 & 5{,}400 & 0 & 34 & 56.7 \\
1.5B & div7 & B & 180 & 0 & 10{,}800 & -- & 47 & 92.5 \\
1.5B & div7 & B & 180 & 1 & 10{,}800 & -- & 44 & 90.0 \\
1.5B & div7 & B & 180 & 2 & 10{,}800 & -- & 41 & 82.5 \\
1.5B & div7 & B & 270 & 0 & 16{,}200 & 0 & 50 & 99.6 \\
1.5B & div7 & B & 270 & 1 & 16{,}200 & 0 & 47 & 92.9 \\
1.5B & div7 & B & 270 & 2 & 16{,}200 & 0 & 45 & 92.1 \\
1.5B & div7 & C & 180 & 0 & 10{,}784 & -- & 72 & 49.6 \\
1.5B & div7 & C & 180 & 1 & 10{,}824 & -- & 30 & 50.4 \\
1.5B & div7 & B$'$ & 180 & 0 & 10{,}800 & 0 & 48 & 48.3 \\
1.5B & div7 & B$'$ & 180 & 1 & 10{,}800 & 0 & 11 & 50.8 \\
1.5B & div7 & D & 180 & 0 & 10{,}800 & 0 & 100 & 50.0 \\
1.5B & div7 & D & 180 & 1 & 10{,}784 & 0 & 0 & 50.4 \\
1.5B & div7 & S & 180 & 0 & 7{,}473 & 1 & 57 & 77.9 \\
1.5B & div7 & S & 180 & 1 & 7{,}215 & 0 & 70 & 65.8 \\
1.5B & div7$\to$div7\_6d & B & 180 & 0 & 0 & 0 & 32 & 77.1 \\
1.5B & div7[cot] & base & 0 & 0 & 0 & 0 & 70 & 67.1 \\
1.5B & div7[cot] & A & 180 & 0 & 360 & 0 & 12 & 55.0 \\
1.5B & div7[cot] & A & 180 & 1 & 360 & 0 & 76 & 53.3 \\
1.5B & div7[fewshot4] & base & 0 & 0 & 0 & 0 & 37 & 47.9 \\
\midrule
1.5B & valid & base & 0 & 0 & 0 & -- & 38 & 30.0 \\
1.5B & valid & A & 60 & 0 & 120 & -- & 50 & 97.5 \\
1.5B & valid & A & 180 & 0 & 360 & -- & 51 & 99.2 \\
1.5B & valid & A & 180 & 1 & 360 & -- & 50 & 100.0 \\
1.5B & valid & A & 180 & 2 & 360 & 0 & 48 & 98.3 \\
1.5B & valid & A & 540 & 0 & 1{,}080 & -- & 50 & 100.0 \\
1.5B & valid & B & 60 & 0 & 2{,}343 & -- & 51 & 87.5 \\
1.5B & valid & B & 180 & 0 & 7{,}001 & -- & 50 & 100.0 \\
1.5B & valid & B & 180 & 1 & 7{,}004 & -- & 55 & 80.0 \\
1.5B & valid & B & 180 & 2 & 6{,}960 & 0 & 37 & 85.0 \\
1.5B & valid & B & 540 & 0 & 21{,}084 & -- & 50 & 100.0 \\
1.5B & valid & C & 180 & 0 & 7{,}092 & -- & 95 & 54.2 \\
1.5B & valid & C & 180 & 1 & 7{,}090 & -- & 1 & 50.8 \\
1.5B & valid & A-full & 1600 & 0 & 3{,}200 & -- & 50 & 100.0 \\
1.5B & valid & S & 180 & 0 & 26{,}233 & 49 & 49 & 64.6 \\
1.5B & valid[cot] & base & 0 & 0 & 0 & 5 & 56 & 66.7 \\
1.5B & valid[fewshot4] & base & 0 & 0 & 0 & 0 & 30 & 68.8 \\
\midrule
1.5B & div2 & base & 0 & 0 & 0 & 0 & 47 & 93.8 \\
1.5B & div2 & A & 180 & 0 & 360 & 0 & 50 & 100.0 \\
1.5B & div2 & A & 180 & 1 & 360 & 0 & 50 & 99.2 \\
1.5B & div2 & B & 180 & 0 & 10{,}800 & 0 & 50 & 100.0 \\
1.5B & div2 & B & 180 & 1 & 10{,}800 & 0 & 50 & 100.0 \\
\midrule
1.5B & div3 & base & 0 & 0 & 0 & 0 & 29 & 78.3 \\
1.5B & div3 & A & 180 & 0 & 360 & 0 & 0 & 50.0 \\
1.5B & div3 & A & 180 & 1 & 360 & 0 & 77 & 62.5 \\
1.5B & div3 & A & 180 & 2 & 360 & 0 & 15 & 49.6 \\
1.5B & div3 & B & 180 & 0 & 10{,}800 & 0 & 50 & 100.0 \\
1.5B & div3 & B & 180 & 1 & 10{,}800 & 0 & 50 & 100.0 \\
1.5B & div3 & B & 180 & 2 & 10{,}800 & 0 & 50 & 100.0 \\
\midrule
1.5B & div11 & base & 0 & 0 & 0 & 32 & 14 & 55.0 \\
1.5B & div11 & A & 180 & 0 & 360 & 0 & 9 & 50.8 \\
1.5B & div11 & A & 180 & 1 & 360 & 0 & 7 & 49.2 \\
1.5B & div11 & A & 180 & 2 & 360 & 0 & 9 & 50.4 \\
1.5B & div11 & B & 180 & 0 & 10{,}800 & 0 & 50 & 99.2 \\
1.5B & div11 & B & 180 & 1 & 10{,}800 & 0 & 49 & 99.2 \\
1.5B & div11 & B & 180 & 2 & 10{,}800 & 0 & 50 & 100.0 \\
\midrule
1.5B & div13 & base & 0 & 0 & 0 & 2 & 23 & 50.4 \\
1.5B & div13 & A & 180 & 0 & 360 & 0 & 0 & 49.6 \\
1.5B & div13 & A & 180 & 1 & 360 & 0 & 58 & 50.4 \\
1.5B & div13 & A & 180 & 2 & 360 & 0 & 8 & 51.2 \\
1.5B & div13 & A & 360 & 0 & 720 & 0 & 0 & 50.0 \\
1.5B & div13 & A & 360 & 1 & 720 & 0 & 100 & 50.0 \\
1.5B & div13 & B & 180 & 0 & 10{,}800 & 0 & 76 & 54.6 \\
1.5B & div13 & B & 180 & 1 & 10{,}800 & 0 & 18 & 49.2 \\
1.5B & div13 & B & 180 & 2 & 10{,}800 & 0 & 49 & 57.9 \\
1.5B & div13 & B & 360 & 0 & 21{,}600 & 0 & 27 & 75.4 \\
1.5B & div13 & B & 360 & 1 & 21{,}600 & 0 & 27 & 72.5 \\
1.5B & div13 & B & 360 & 2 & 21{,}600 & 0 & 28 & 72.5 \\
1.5B & div13 & B & 540 & 0 & 32{,}400 & 0 & 46 & 94.6 \\
1.5B & div13 & B & 540 & 1 & 32{,}400 & 0 & 45 & 95.0 \\
1.5B & div13 & B & 540 & 2 & 32{,}400 & 0 & 46 & 95.8 \\
1.5B & div13 & C & 180 & 0 & 10{,}824 & 0 & 49 & 45.0 \\
1.5B & div13 & C & 180 & 1 & 10{,}788 & 5 & 86 & 47.9 \\
1.5B & div13[cot] & base & 0 & 0 & 0 & 0 & 15 & 54.6 \\
1.5B & div13[cot] & A & 180 & 0 & 0 & 0 & 0 & 50.0 \\
\midrule
1.5B & div7\_6d & base & 0 & 0 & 0 & 0 & 88 & 52.5 \\
1.5B & div7\_6d & A & 180 & 0 & 360 & 0 & 7 & 47.9 \\
1.5B & div7\_6d & A & 180 & 1 & 360 & 0 & 9 & 50.8 \\
1.5B & div7\_6d & B & 180 & 0 & 14{,}040 & 0 & 48 & 96.2 \\
1.5B & div7\_6d & B & 180 & 1 & 14{,}040 & 0 & 43 & 90.0 \\
\midrule
3B & prime & base & 0 & 0 & 0 & 110 & 2 & 47.1 \\
3B & prime & A & 180 & 0 & 360 & 0 & 67 & 82.9 \\
3B & prime & A & 180 & 1 & 360 & 0 & 63 & 81.7 \\
3B & prime & B & 180 & 0 & 13{,}370 & 0 & 60 & 90.4 \\
3B & prime & B & 180 & 1 & 13{,}260 & 0 & 41 & 79.6 \\
3B & prime & S & 180 & 0 & 19{,}736 & 49 & 22 & 69.6 \\
3B & prime$\to$prime\_hard & A & 180 & 0 & 0 & 0 & 99 & 50.4 \\
3B & prime$\to$prime\_hard & A & 180 & 1 & 0 & 0 & 92 & 47.9 \\
3B & prime$\to$prime\_hard & B & 180 & 0 & 0 & 0 & 100 & 50.4 \\
3B & prime$\to$prime\_hard & B & 180 & 1 & 0 & 0 & 69 & 53.8 \\
\midrule
3B & div7 & base & 0 & 0 & 0 & 0 & 55 & 94.2 \\
3B & div7 & A & 180 & 0 & 360 & 0 & 5 & 54.6 \\
3B & div7 & A & 180 & 1 & 360 & 0 & 100 & 50.0 \\
3B & div7 & B & 90 & 0 & 5{,}400 & 0 & 45 & 60.0 \\
3B & div7 & B & 90 & 1 & 5{,}400 & 0 & 67 & 58.8 \\
3B & div7 & B & 180 & 0 & 10{,}800 & 0 & 27 & 68.3 \\
3B & div7 & B & 180 & 1 & 10{,}800 & 0 & 41 & 82.1 \\
3B & div7 & B & 270 & 0 & 16{,}200 & 0 & 48 & 97.9 \\
3B & div7 & B & 270 & 1 & 16{,}200 & 0 & 47 & 86.7 \\
3B & div7 & B$'$ & 180 & 0 & 10{,}800 & 51 & 96 & 50.8 \\
3B & div7 & B$'$ & 180 & 1 & 10{,}800 & 0 & 45 & 48.3 \\
3B & div7 & D & 180 & 0 & 10{,}800 & 0 & 100 & 50.0 \\
3B & div7 & D & 180 & 1 & 10{,}784 & 0 & 100 & 50.0 \\
3B & div7 & S & 180 & 0 & 11{,}130 & 0 & 52 & 97.5 \\
3B & div7 & S & 180 & 1 & 11{,}198 & 0 & 54 & 96.2 \\
3B & div7$\to$div7\_6d & A & 180 & 0 & 0 & 0 & 0 & 50.4 \\
3B & div7$\to$div7\_6d & B & 180 & 0 & 0 & 0 & 22 & 54.2 \\
3B & div7$\to$div7\_6d & B & 180 & 1 & 0 & 0 & 33 & 65.0 \\
3B & div7[cot] & A & 180 & 0 & 360 & 0 & 3 & 53.3 \\
\midrule
3B & valid & base & 0 & 0 & 0 & 10 & 52 & 85.8 \\
3B & valid & A & 180 & 0 & 360 & 0 & 50 & 100.0 \\
3B & valid & A & 180 & 1 & 360 & 0 & 50 & 100.0 \\
3B & valid & B & 180 & 0 & 7{,}001 & 0 & 59 & 91.2 \\
3B & valid & B & 180 & 1 & 7{,}004 & 0 & 50 & 100.0 \\
3B & valid & S & 180 & 0 & 30{,}284 & 0 & 57 & 92.1 \\
\midrule
3B & div3 & A & 180 & 0 & 360 & 0 & 50 & 99.2 \\
3B & div3 & A & 180 & 1 & 360 & 0 & 50 & 99.2 \\
3B & div3 & B & 180 & 0 & 10{,}800 & 0 & 50 & 100.0 \\
3B & div3 & B & 180 & 1 & 10{,}800 & 0 & 50 & 100.0 \\
\midrule
3B & div11 & A & 180 & 0 & 360 & 0 & 6 & 55.4 \\
3B & div11 & A & 180 & 1 & 360 & 0 & 1 & 50.8 \\
3B & div11 & B & 180 & 0 & 10{,}800 & 0 & 50 & 99.6 \\
3B & div11 & B & 180 & 1 & 10{,}800 & 0 & 50 & 100.0 \\
\midrule
3B & div13 & base & 0 & 0 & 0 & 0 & 53 & 92.1 \\
3B & div13 & A & 180 & 0 & 360 & 0 & 0 & 50.0 \\
3B & div13 & A & 180 & 1 & 360 & 0 & 28 & 48.3 \\
3B & div13 & A & 360 & 0 & 720 & 0 & 0 & 50.0 \\
3B & div13 & A & 540 & 0 & 1{,}080 & 0 & 0 & 50.0 \\
3B & div13 & B & 180 & 0 & 10{,}800 & 0 & 70 & 52.9 \\
3B & div13 & B & 180 & 1 & 10{,}800 & 0 & 42 & 59.2 \\
3B & div13 & B & 360 & 0 & 21{,}600 & 0 & 33 & 83.3 \\
3B & div13 & B & 360 & 1 & 21{,}600 & 0 & 45 & 92.9 \\
3B & div13 & B & 540 & 0 & 32{,}400 & 0 & 50 & 99.6 \\
3B & div13 & B & 540 & 1 & 32{,}400 & 0 & 49 & 99.2 \\
3B & div13[cot] & A & 180 & 0 & 360 & 0 & 0 & 50.0 \\
\midrule
3B & div7\_6d & A & 180 & 0 & 360 & 0 & 1 & 49.2 \\
3B & div7\_6d & B & 180 & 0 & 14{,}040 & 0 & 35 & 82.5 \\
\midrule
7B & prime & base & 0 & 0 & 0 & 4 & 30 & 72.1 \\
7B & prime & A & 180 & 0 & 360 & 0 & 1 & 51.2 \\
7B & prime & A & 180 & 1 & 360 & 0 & 67 & 82.5 \\
7B & prime & B & 180 & 0 & 13{,}370 & 1 & 60 & 89.2 \\
7B & prime & B & 180 & 1 & 13{,}260 & 0 & 36 & 80.0 \\
7B & prime & S & 180 & 0 & 26{,}737 & 25 & 32 & 72.1 \\
\midrule
7B & div7 & base & 0 & 0 & 0 & 1 & 45 & 82.5 \\
7B & div7 & A & 180 & 0 & 360 & 0 & 35 & 49.2 \\
7B & div7 & A & 180 & 1 & 360 & 0 & 75 & 52.9 \\
7B & div7 & B & 180 & 0 & 10{,}800 & 0 & 43 & 86.7 \\
7B & div7 & B & 180 & 1 & 10{,}800 & 0 & 46 & 90.4 \\
7B & div7 & B$'$ & 180 & 0 & 10{,}800 & 3 & 16 & 53.3 \\
7B & div7 & B$'$ & 180 & 1 & 10{,}800 & 0 & 54 & 46.7 \\
7B & div7 & D & 180 & 0 & 10{,}800 & 0 & 39 & 47.9 \\
7B & div7 & D & 180 & 1 & 10{,}784 & 0 & 61 & 50.4 \\
7B & div7 & S & 180 & 0 & 10{,}803 & 0 & 41 & 82.1 \\
7B & div7 & S & 180 & 1 & 10{,}728 & 0 & 42 & 81.2 \\
7B & div7[cot] & A & 180 & 0 & 360 & 0 & 60 & 55.0 \\
\midrule
7B & valid & base & 0 & 0 & 0 & 3 & 42 & 89.6 \\
7B & valid & A & 180 & 0 & 360 & 0 & 50 & 100.0 \\
7B & valid & A & 180 & 1 & 360 & 0 & 50 & 100.0 \\
7B & valid & B & 180 & 0 & 7{,}001 & 0 & 92 & 58.3 \\
7B & valid & B & 180 & 1 & 7{,}004 & 0 & 52 & 97.5 \\
\midrule
7B & div11 & A & 180 & 0 & 360 & 0 & 26 & 55.4 \\
7B & div11 & A & 180 & 1 & 360 & 0 & 16 & 51.7 \\
7B & div11 & B & 180 & 0 & 10{,}800 & 0 & 49 & 98.8 \\
7B & div11 & B & 180 & 1 & 10{,}800 & 0 & 50 & 100.0 \\
\midrule
7B & div13 & base & 0 & 0 & 0 & 0 & 47 & 80.0 \\
7B & div13 & A & 180 & 0 & 360 & 0 & 25 & 45.4 \\
7B & div13 & A & 180 & 1 & 360 & 0 & 6 & 47.9 \\
7B & div13 & B & 180 & 0 & 10{,}800 & 0 & 26 & 61.3 \\
7B & div13 & B & 180 & 1 & 10{,}800 & 0 & 35 & 65.8 \\
7B & div13[cot] & A & 180 & 0 & 360 & 0 & 68 & 46.2 \\
\midrule
SmolLM2 & prime & A & 180 & 0 & 360 & 0 & 36 & 68.3 \\
SmolLM2 & prime & B & 180 & 0 & 13{,}370 & 0 & 23 & 62.5 \\
\midrule
SmolLM2 & div7 & base & 0 & 0 & 0 & 0 & 4 & 49.6 \\
SmolLM2 & div7 & A & 180 & 0 & 360 & 0 & 29 & 45.0 \\
SmolLM2 & div7 & A & 180 & 1 & 360 & 0 & 91 & 47.1 \\
SmolLM2 & div7 & B & 180 & 0 & 10{,}800 & 0 & 41 & 60.0 \\
SmolLM2 & div7 & B & 180 & 1 & 10{,}800 & 0 & 51 & 69.6 \\
SmolLM2 & div7 & D & 180 & 0 & 10{,}800 & 0 & 87 & 51.7 \\
\midrule
SmolLM2 & valid & A & 180 & 0 & 360 & 0 & 50 & 99.6 \\
SmolLM2 & valid & B & 180 & 0 & 7{,}001 & 0 & 30 & 80.4 \\
\end{longtable}
}{\fbox{Run table missing (run \texttt{paper/exp/make\_numbers.py}).}}}

\begin{table}[h]
\caption{Timestamped predictions (S1--S6 before any of their runs; S7--S12 appended during the study, each before its cells) and their status. 3B and 7B were not the preregistered model; their columns report the same tests at those sizes; [$k$]: number of seeds. \nr: not run.}
\label{tab:prereg}
\centering\scriptsize\setlength{\tabcolsep}{2.5pt}
\begin{tabular}{@{}lp{0.24\linewidth}p{0.26\linewidth}lp{0.14\linewidth}p{0.12\linewidth}@{}}
\toprule
& prediction & result (1.5B) & status & 3B & 7B \\
\midrule
S1 & div7 B $\ge 80$, B$-$A $\ge 20$\pp{} (seeds 1, 2) & B \divsevenBsOne/\divsevenBsTwo, A \divsevenAsOne/\divsevenAsTwo & \verdictSOne & & \\
S2 & traces help $\ge 2$ of div3/11/13 (2 seeds) & B$-$A div3 \divthreeBAdZero/\divthreeBAdOne, div11 \divelevenBAdZero/\divelevenBAdOne, div13 \divthirteenBAdZero/\divthirteenBAdOne & \verdictSTwo & \verdictSTwoThreeB{} (div3 \threeBdivthreeBootBADiff) & div3 \nr \\
S3 & div2 A $\ge 95$ & \divtwoA & \verdictSThree{} (settled) & & \\
S4 & div7 B $>$ A at 3B (2 seeds) & (3B) A \threeBdivsevenAsZero/\threeBdivsevenAsOne, B \threeBdivsevenBsZero/\threeBdivsevenBsOne & \verdictSFour & & \\
S5 & 6-digit div7 B within 10\pp{} of $q^6=\divsevenSixPred$ & \divsevenSixBsZero/\divsevenSixBsOne & \textbf{\verdictSFive} & \verdictSFiveThreeB & \nr \\
S6 & prime A on hard set $<60$, follows rule & \primeHardA, rule agreement \primeHardAagree & \verdictSSix{} (settled) & \verdictSSixThreeB & \nr \\
S7 & $q$ follows $m$; $\ge 4$ of 5 cells & (a)--(c) hold, (e) at ceiling, (d) \SsevenCellD: 3/4 informative & \verdictSSeven{} (\SsevenHeld/5) & post hoc only & \\
S8 & arm D $\le 60$ (2 seeds) & \divsevenDsZero/\divsevenDsOne & \verdictSEight & D \threeBdivsevenD{} [\threeBdivsevenDnSeeds] & D \sevenBdivsevenD{} [\sevenBdivsevenDnSeeds] \\
S9 & div13, $n=360$: B $\ge 80$, A $\le 60$ & B \divthirteenBnThreeSixtysZero, A \divthirteenAnThreeSixtysZero & \textbf{\verdictSNine} & \verdictSNineThreeB{} (\threeBdivthirteenBnThreeSixtysZero) & \nr \\
S10 & arm S div7 $\le 65$ (2 seeds); valid $\ge 90$ & div7 \divsevenSelfTrace{} (\divsevenSelfSeeds) & \textbf{\verdictSTen} & contradicted (div7 \threeBdivsevenSsZero/\threeBdivsevenSsOne) & div7 \sevenBdivsevenS{} [\sevenBdivsevenSnSeeds] (base \sevenBdivsevenBase) \\
S11 & second family: div7 B$-$A $\ge 20$\pp{} (2 seeds), D within 10\pp{} of A & \secondFamilyModel: B$-$A \famDivsevenBAdiffsZero/\famDivsevenBAdiffsOne, D \famDivsevenDsZero{} vs A \famDivsevenAsZero & \textbf{\SElevenVerdict} & & \\
S12 & div7 A sweep (lr, epochs; 2 seeds): no cell $\ge 65$, one non-constant, div3 A $\ge 90$ & max \sweepAmax{} (\sweepAmaxConfig); non-constant \sweepAnonCollapse/4; div3 \sweepDivthreeA & \STwelveVerdict & & \\
\bottomrule
\end{tabular}
\end{table}

\begin{table}[h]
\caption{Timestamped cells in detail. Qwen2.5-1.5B-Instruct, 4 digits and $n=\nTrain$ unless stated; accuracy in \%; seed 0 unless stated. For S7 the B column is the per-step accuracy $q$ (written $p$ in the preregistration text); for S8 and S10 it is the accuracy of arms D and S. \nr: not run.}
\label{tab:prereg-detail}
\centering\footnotesize\setlength{\tabcolsep}{3pt}
\begin{tabular}{@{}llcccp{0.27\linewidth}c@{}}
\toprule
& cell & probe & A & B & criterion & met \\
\midrule
S1 & div7, s1 / s2 & \divsevenProbe & \divsevenAsOne{} / \divsevenAsTwo & \divsevenBsOne{} / \divsevenBsTwo & B $\ge 80$, B$-$A $\ge 20$\pp, $p<0.01$ in each seed & \verdictSOne \\
S2 & div3 & \divthreeProbe & \divthreeA & \divthreeB & \multirow{3}{=}{$\ge 2$ of the 3 with B$-$A $\ge 15$\pp{} in both seeds} & \multirow{3}{*}{\verdictSTwo} \\
   & div11 & \divelevenProbe & \divelevenA & \divelevenB & & \\
   & div13 & \divthirteenProbe & \divthirteenA & \divthirteenB & & \\
S3 & div2 & \divtwoProbe & \divtwoA & \divtwoB & A $\ge 95$ & \verdictSThree \\
S4 & div7, 3B model & & \threeBdivsevenA & \threeBdivsevenB & B $>$ A & \verdictSFour \\
S5 & div7 (6 digits), s0 / s1 & & \divsevenSixA & \divsevenSixBsZero{} / \divsevenSixBsOne & $|\mathrm{B}-q^6|\le 10$\pp; $q^6=\divsevenSixPred$ & \verdictSFive \\
   & secondary: 4-digit B & & & \divsevenTransferSixB & same & \\
S6 & prime$\to$hard set & & \primeHardA & \nr & A $<60$; agrees more with the rule than with the label & \verdictSSix \\
S7 & (a) div13, $n=540$ & & & \pDivthirteennFiveFortysZero & $q\ge 0.95$, acc.\ $\ge 85$ (\divthirteenBnFiveFortysZero) & \SsevenCellA \\
   & (b) div7, $n=270$ & & & \pDivsevennTwoSeventysZero & $q\ge 0.97$ & \SsevenCellB \\
   & (c) div7, $n=90$ & & & \pDivsevennNinetysZero & $q\le 0.80$, acc.\ $\le 75$ (\divsevenBnNinetysZero) & \SsevenCellC \\
   & (d) div11 & & & \pDiveleven & $0.566<q<0.955$ & \SsevenCellD \\
   & (e) div3 / div2 & & & \pDivthree{} / \pDivtwo & $q\ge 0.97$ & \SsevenCellE \\
S8 & div7, arm D, s0 / s1 & & & \divsevenDsZero{} / \divsevenDsOne & D $\le 60$ in both seeds & \verdictSEight \\
S9 & div13, $n=360$ & & \divthirteenAnThreeSixtysZero & \divthirteenBnThreeSixtysZero & B $\ge 80$, A $\le 60$ & \verdictSNine \\
S10 & div7, arm S & & & \divsevenSelfTrace & S $\le 65$ in seeds 0 and 1; valid S $\ge 90$ & \verdictSTen \\
\bottomrule
\end{tabular}
\end{table}

\begin{table}[h]
\caption{Earlier answer-only study (same model and LoRA configuration; \priorTrainN{} training items, answer-only prompts, 3--4-digit numbers for arithmetic; mean of two seeds). Correct labels vs.\ labels permuted across training items.}
\label{tab:shuffled}
\centering\small
\begin{tabular}{lccc}
\toprule
task & base & correct labels & shuffled labels \\
\midrule
divisible by 3 & \priorDivthreeBase & \priorDivthreeA & \priorDivthreeShuf \\
divisible by 7 & \priorDivsevenBase & \priorDivsevenA & \priorDivsevenShuf \\
divisible by 13 & \priorDivthirteenBase & \priorDivthirteenA & \priorDivthirteenShuf \\
perfect square & \priorSquareBase & \priorSquareA & \priorSquareShuf \\
prime & \priorPrimeBase & \priorPrimeA & \priorPrimeShuf \\
argument validity & \priorValidBase & \priorValidA & \priorValidShuf \\
joint consistency (SAT) & \priorConsistBase & \priorConsistA & \priorConsistShuf \\
\bottomrule
\end{tabular}
\end{table}

\section{Setup details}
\label{app:setup}
\begin{table}[h]
\caption{Training and evaluation setup, shared by every arm and model.}
\label{tab:setup}
\centering\footnotesize
\begin{tabular}{@{}lp{0.68\linewidth}@{}}
\toprule
models & Qwen2.5-0.5B/1.5B/3B/7B-Instruct, bf16; 0.5B and 1.5B on a laptop GPU, 3B and 7B on one RTX 5090 \\
adapter & LoRA $r=\loraRank$, $\alpha=\loraAlpha$, dropout \loraDropout, on q/k/v/o projections only \\
optimization & AdamW, learning rate $\learningRate$, batch size 1, \nEpochs{} epochs, loss on completion tokens only \\
prompt & the model's chat template; one user turn ending ``End your reply with `Answer: Yes' or `Answer: No'.'' \\
seed & selects the balanced training sample and seeds PyTorch (adapter initialization, data order) \\
decoding & greedy; generation cap \maxNewDiv{} tokens (div2/3/7), \maxNewDivthirteen{} (div11/13), \maxNewDivsevenSix{} (6-digit div7), \maxNewPrime{} (prime), \maxNewValid{} (valid), 512 for the CoT prompt \\
extractor & last ``Answer: Yes/No'' in the completion, else the first standalone Yes/No; otherwise unparsed; unparsed and truncated answers count as wrong \\
pools & valid: \validPool{} training arguments; prime: \primeAtokN{} numbers (arm A-full); test items never in a pool \\
arm D & for each training item, the correct trace of another pool number with the same label \\
arm S & $K=4$ samples per item at temperature 0.7 from the base model's CoT prompt; one accepted completion per item \\
\bottomrule
\end{tabular}
\end{table}
Base-model cells whose generations often hit the cap are lower bounds: 3B prime base \threeBprimeBase\% (\threeBprimeBaseCap{} of \nTest{} capped), 1.5B prime with \nShots{} demonstrations \primeFew\% (\primeFewCap{} capped) and 1.5B div11 base \divelevenBase\% (\divelevenBaseCap{} capped). The 1.5B valid base cell (\validBase\%) has no saved generations and is unaudited. Arm A collapses to a constant answer in single seeds (Yes rate near 0 or 100\% in Table~\ref{tab:runs}): 1.5B div3 and div13 seed 0, 3B div7 seed 1 and div13 seed 0; at larger $n$, 1.5B div7 at $n=540$, 1.5B div13 at $n=360$ and 3B div13 at $n=360$ and 540 (all ``No''); 3B div7$\to$div7\_6d (``No'') and 3B div7\_6d (Yes rate 1\%); and 7B prime seed 0, which answers ``No'' on \sevenBprimeAsZeroNo{} of \nTest{} items. A single untuned configuration is used throughout, so these collapses are part of what answer-only training did here, not a tuned optimum.

\section{The logic task}
\label{app:logic}
Arguments instantiate \nFormsHalf{} valid forms (modus ponens, modus tollens, hypothetical syllogism, disjunctive syllogism, constructive dilemma, contraposition) and \nFormsHalf{} invalid ones (affirming the consequent, denying the antecedent, affirming a disjunct, a broken chain, the converse, and ``or'' to ``and'') with distinct random atom letters, rendered in English (``if $A$ then $B$'', ``not $A$'', ``$A$ or $B$''). The schema proposes an argument and Z3 decides it: the argument is labeled valid iff the premises together with the negated conclusion are unsatisfiable. For a valid argument the trace names the unsatisfiable core; for an invalid one it gives the satisfying assignment as a countermodel. Training and test share the \nForms{} forms and never an instance, so a model can reach high accuracy by recognizing the form, which the atom-renamed $n$-gram probe does (\validProbe\%). The task therefore checks that answers work when the surface determines the label; it is not a test of logical reasoning. Its results: base \validBase\%, A \validA\% (\validAsZero{} and \validAsOne), B \validB\% (\validBsZero{} and \validBsOne), C \validC\% (\validCsZero{} and \validCsOne); at $n=\nTrainSmall$, A \validAsmall\% and B \validBsmall\%; at $n=\nTrainLarge$, both \validAlarge\%.

\section{Trace formats}
\label{app:formats}
All arms share the prompt, which ends ``End your reply with `Answer: Yes' or `Answer: No'.'' Examples of training completions, verbatim:

\begin{small}
\noindent\textbf{div7, B} (prompt: ``Is 1617 divisible by 7?''):
\begin{quote}\ttfamily\raggedright
Process the digits of 1617 keeping a running remainder mod 7. r = (10*0 + 1) mod 7 = 1. r = (10*1 + 6) mod 7 = 2. r = (10*2 + 1) mod 7 = 0. r = (10*0 + 7) mod 7 = 0. Final remainder 0, so 1617 is divisible by 7.\\Answer: Yes
\end{quote}
\noindent\textbf{div7, C} (a scrambled trace from another example; the answer line keeps this example's label):
\begin{quote}\ttfamily\raggedright
running mod r = so = 4067 Process by 7. 7 a = 0. 0) 0. 7 mod (10*5 = r 7 the = = keeping + r is mod mod 4) (10*0 + of 4067 r + 7) 5. = + digits remainder 7. (10*0 = 7 0, divisible Final 6) 4. (10*4 mod remainder\\Answer: Yes
\end{quote}
\noindent\textbf{prime, B}, a composite (``Is 4143 a prime number?''):
\begin{quote}\ttfamily\raggedright
Check prime divisors up to isqrt(4143) = 64. 4143 mod 2 = 1. 4143 mod 3 = 0. 3 divides 4143, so 4143 is composite.\\Answer: No
\end{quote}
\noindent\textbf{prime, B}, a prime (``Is 1493 a prime number?''):
\begin{quote}\ttfamily\raggedright
Check prime divisors up to isqrt(1493) = 38. 1493 mod 2 = 1. 1493 mod 3 = 2. 1493 mod 5 = 3. 1493 mod 7 = 2. 1493 mod 11 = 8. 1493 mod 13 = 11. 1493 mod 17 = 14. 1493 mod 19 = 11. 1493 mod 23 = 21. 1493 mod 29 = 14. 1493 mod 31 = 5. 1493 mod 37 = 13. No prime divisor up to 38, so 1493 is prime.\\Answer: Yes
\end{quote}
\noindent\textbf{valid, B}, invalid (``Premises: (1) if H then D; (2) D. Conclusion: H.''):
\begin{quote}\ttfamily\raggedright
Assume every premise is true and the conclusion false. That is satisfiable: set D=true, H=false. Under this assignment all premises hold while the conclusion is false, a countermodel, so the argument is invalid.\\Answer: No
\end{quote}
\noindent\textbf{valid, B}, valid (``Premises: (1) if S then R; (2) S. Conclusion: R.''):
\begin{quote}\ttfamily\raggedright
Assume every premise is true and the conclusion false, i.e. \textasciitilde(R) holds. Together these force a contradiction: \{S -> R, S, \textasciitilde(R)\} cannot all hold at once. Since assuming the conclusion false is impossible, the conclusion follows.\\Answer: Yes
\end{quote}
\noindent\textbf{A} (every task): \texttt{Answer: Yes} or \texttt{Answer: No}.
\end{small}

\section{Preregistration text and decision log}
\label{app:prereg}
\paragraph{Provenance of the prior div7 result.} A verdict written on \priorVerdictDate{} from the \nTrainSmall-example trace run alone (\divsevenBsmall\%) stated that traces do not unlock divisibility by 7, while the preregistered \nTrain-example cell was still running; that cell finished at \divsevenBsZero\% with no code change. The dose curve shows why both were right about their own $n$. The replication of that cell was the go/no-go condition S1, and it passed (\divsevenBsOne\% and \divsevenBsTwo\% in seeds 1 and 2).

\paragraph{First preregistration (prior grid).} Written before any run of the prior grid, according to the history of the repository in which it was written, and never edited afterwards; results were appended below it. The file enters the present repository with the first commit of the manifest (P1), so its earlier date rests on the original repository. It fixed: arms A, B, C, an answers-only arm at larger scale, arms using model-generated traces, and a conditional wrong-trace arm; Qwen2.5-1.5B-Instruct with the LoRA configuration of \S\ref{sec:setup}; tasks prime, valid and div7, all reported; \nTest{} problem-disjoint test items per task (\nTestClass{} per class), identical across arms, graded by one parser on the final answer. Endpoints: H1, B $>$ A at $n=\nTrain$ per task, paired McNemar, Holm across the three tasks; H2, B $>$ C at $n=\nTrain$; H3, accuracy at matched training tokens; H4, on div7 B reaches about 60\% with A at chance. The full text is in the supplementary material. Outcomes: H1 held on div7 only, H2 on div7 and valid, H3 failed where answers work (not tested on div7), and H4 held at $n=\nTrain$ only.

\paragraph{Second preregistration and its decision log.} Reproduced from \texttt{PREDICTIONS.md} up to the S10 entry; the six later entries (the S10 primary outcome, a wording note, S11, a second-model-family test none of whose cells had run at submission, the post hoc correction of the S10 trace-audit numbers, and the two entries recording the remaining S10 cells) are summarized in \S\ref{sec:controls} and reproduced in the supplementary material. Markdown markup converted, commit hashes replaced by the anonymized identifiers P1--P6 of the timestamp manifest and the location in each time stamp removed, text otherwise unchanged. The file's $p$ (per-step accuracy) is $q$ in the paper. The clock times in the header and the first entries were first written as estimates (the header as ``\textasciitilde20:35'') and later replaced by commit timestamps; the decision log notes this. After the runs, a dated note was appended recording that the seed-0 ``re-run'' was a deterministic re-execution.
\begin{small}
\begin{quote}
\textbf{Preregistered predictions and decision rules (frozen before any run listed here).} Committed 2026-09-25 20:29:47 (commit P1), before any of the runs below were started. The commit hash of this file is the preregistration. Earlier results (results/THINKING\_VS\_DATA.md, frozen 2026-08-04, and its later staged results block) are prior data, reported as such.

\textbf{Provenance of the prior div7 result (stated in the paper).} The 2026-08-04 verdict ``H4 (a procedure trace unlocks div7): NO'' was written from the n = 60 trace run alone (B@60 = 50.0\%). The preregistered main cell, B@180 seed 0, finished afterwards at 92.5\% with the same worker (no code change). The paper reports the premature verdict and the later result.

\textbf{Claim under test.} Exact labels teach shortcuts; exact traces teach short programs. With labels from an exact verifier:
answer-only supervision (arm A) recovers roughly what a simple probe on the answer surface recovers;
trace supervision (arm B: procedure trace, then answer) succeeds when the per-step trace accuracy p, compounded over k steps, stays high (predicted accuracy about p\textasciicircum k), and fails when the procedure is long.

\textbf{Decision rules.}
\textbf{S1 (go/no-go, before any other new run):} div7 B@180 seeds 1 and 2 each reach >= 80\% accuracy, with B - A >= 20 percentage points and McNemar p < 0.01 against A of the same seed. If S1 fails, the paper is not submitted in this form.
\textbf{S2:} at least 2 of \{div3, div11, div13\} (4-digit numbers) give B - A >= 15 pp with the same sign in both seeds.
\textbf{S3:} div2 A >= 95\% (the answer is surface-learnable from the last digit).
\textbf{S4:} the div7 sign (B > A) holds for Qwen2.5-3B-Instruct.
\textbf{S5 (falsifiable length test):} on 6-digit div7, B's accuracy is within 10 pp of p\textasciicircum 6, where p is B's per-step accuracy measured on 4-digit div7 generations.
\textbf{S6 (shortcut):} A on prime agrees more with the rule ``odd and no factor <= 7'' than with the true label on hard negatives (odd composites with no factor <= 7), where A's accuracy drops below 60\%.

Definitions: a ``win'' is a difference of >= 10 pp with both seeds agreeing in sign; ``neither'' is both arms at or below majority class + 8 pp. McNemar tests per seed, Holm correction across tasks within each hypothesis family, alpha = 0.05. Every cell is reported, including those that go against the claim.

\textbf{Arms.} base (zero-shot), A (answers), B (trace then answer), C (length-matched scrambled trace), B' (answer then trace, div7 only), plus eval-only controls: zero-shot chain-of-thought prompt and 4-shot trace demonstrations on the base model. Model: Qwen2.5-1.5B-Instruct (laptop), Qwen2.5-3B-Instruct and others where the RTX 5090 allows. LoRA r = 8, alpha 16, 3 epochs, lr 2e-4, bf16, greedy decoding, n = 180 training examples, 240 problem-disjoint test items.

\textbf{Decision log.}

2026-09-25 21:09:47 (commit P2) --- \textbf{S1 PASS.} div7 B@180 seed 1 = 90.0\% vs A 50.8\% (+39.2 pp, McNemar p = 2e-18); seed 2 = 82.5\% vs A 47.9\% (+34.6 pp, p = 6.8e-13). Runs in results/thinking\_vs\_data/runs\_L1.jsonl, produced by the unchanged worker experiments/exp\_thinking\_ft\_worker.py. The paper proceeds; the remaining tests run on experiments/exp\_worker\_v2.py.

2026-09-25 21:26:10 (commit P3; the div7\_6d B seed-0 generation file was written at 21:27:30) --- \textbf{S5 clarification, recorded before either S5 cell finished} (div7\_6d B seed 0 was mid-training; the 4-digit-adapter transfer eval had not started). Primary S5 cell = arm B trained and tested on the 6-digit task (div7\_6d), seeds 0 and 1, compared with p\textasciicircum 6 where p = per-step accuracy of 4-digit div7 B generations (seed 0: p = 0.955, p\textasciicircum 6 = 0.758). Secondary = the 4-digit B adapter evaluated on 6-digit inputs. Also reported, labelled post hoc: the fraction of fully correct 6-digit traces vs p\textasciicircum 6, and answer accuracy vs p\textasciicircum 6 + (1 - p\textasciicircum 6) * g, where g is the answer-correct rate when the trace is wrong (4-digit: g = 0.561), since a wrong trace still yields the right yes/no answer about half the time.

2026-09-25 21:29:52 (commit P4) --- \textbf{S5 (primary) FAILED as preregistered}: 6-digit div7 B seed 0 = 96.2\% vs p\textasciicircum 6 = 75.8\% (+20.4 pp > 10 pp tolerance). Its own per-step accuracy is 98.8\% (vs 95.5\% for the 4-digit model), so p is not a fixed property that transfers between training sets.

2026-09-25 21:29:52 (commit P4; div11 B seed 0 started after 21:30:58) --- \textbf{S7 (new hypothesis, preregistered before any of its cells ran):} per-step trace accuracy p is set by the number of supervised transitions per entry of the (remainder, digit) table, m = k * n / (10 * d) (k steps per trace, n training examples, 10d table entries). Observed so far: div7 n=180 m=10.3 -> p=0.955; div7\_6d n=180 m=15.4 -> p=0.988; div13 n=180 m=5.5 -> p=0.566. Predictions: (a) div13 B n=540 seed 0 (m=16.6): p >= 0.95 and answer accuracy >= 85\%; (b) div7 B n=270 seed 0 (4-digit, m=15.4): p >= 0.97; (c) div7 B n=90 seed 0 (m=5.1): p <= 0.80 and answer accuracy <= 75\%; (d) div11 B n=180 (m=6.5, not yet run): 0.566 < p < 0.955; (e) div3 B n=180 (m=24) and div2 B n=180 (m=36): p >= 0.97. S7 passes if at least 4 of (a)-(e) hold; every cell is reported either way.

Note (21:35): the clock times in the entries above were first written as estimates; they have been replaced by the commit timestamps, which are the authoritative record.

2026-09-25 21:38:05 (commit P5) \textbf{S8 and S9, preregistered before their cells ran} (div11 B seed 0 = 99.2\% had finished; none of the cells below had started): S8 (grounding control): div7 arm D (a fluent, correct long-division trace of a DIFFERENT number with the same label) at n = 180, seeds 0 and 1: accuracy <= 60\% in both seeds, i.e. within 10 pp of A, far below B. S9 (dose, step difficulty fixed): div13 B at n = 360, seed 0 (m = 11.1 transitions per table entry, close to div7's 10.3 at n = 180): accuracy >= 80\%; div13 A at n = 360 seed 0 stays <= 60\%.

2026-09-25 21:46:44 (commit P6) \textbf{S10, preregistered before any arm-S cell ran.} Arm S = self-generated traces filtered by the exact verifier on the final answer only (STaR / rejection-sampling fine-tuning, K = 4 samples per item at temperature 0.7 from the base model's chain-of-thought prompt; one kept completion per item; same 180 items and LoRA as A/B). Prediction: outcome-only verification does not supply correct steps, so on div7 arm S does not unlock the task: div7 S accuracy <= 65\% in seeds 0 and 1, and fewer than half of the kept div7 traces have all remainders correct. On valid, where answers are surface-learnable, S >= 90\%. Prime S is reported without a prediction.
\end{quote}
\end{small}

\section{An automaton view (intuition)}
\label{app:automaton}
\begin{remark}[intuition, not a theorem]
\label{rem:automaton}
Let $x_1\cdots x_k$ be the digits of $n$ and $\delta_d(r,a)=(10r+a)\bmod d$. Then $d\mid n$ iff $\delta_d^{*}(0,x_1\cdots x_k)=0$: divisibility by $d$ is recognized by a $d$-state automaton reading digits, whose states the arm-B trace writes. Under teacher forcing each supervised state is a function of two tokens already in the context, so $n$ traces are $kn$ samples of one table with $10d$ entries, $m=kn/(10d)$ per entry. Answer-only supervision gives $n$ samples of the $k$-fold composition, and for $d$ coprime to 10 the label is nearly independent of every proper subset of the digits: for $d=7$ and $k=\divsevenSteps$, fixing any three digits leaves $\Pr(7\mid n)$ between \condLow{} and \condHigh{} (marginal \condMarg).
\end{remark}
The remark resembles the globality barrier of \citet{abbe2024far} and the teacher-forced CoT of \citet{kim2025parity}, but those concern growing input length and do not apply literally to fixed 4-digit inputs. It counts how often each transition is supervised, not how hard it is: for $d=3$ and $d=11$ the transition is an addition or subtraction mod $d$, for $d=7$ and $d=13$ it involves a multiplication, and the results follow difficulty as well as $m$. Its independence argument for answer-only supervision applies to $d=3$ as much as to $d=7$, yet at 3B answers alone teach div3 (\threeBdivthreeA\%). The count is the wrong measure there: divisibility by 3 depends on the digits only through their sum, a symmetric, low-degree function, while divisibility by 7 or 13 does not, and the locality and degree accounts of \citet{abbe2024far} and \citet{prystawski2023why} separate the two where the independence count does not. The remark is intuition, not a theorem.

\section{RLVR: a hypothesis, not a result}
\label{app:rlvr}
With a binary verifier the RLVR reward depends on the final answer only, so we hypothesize, without testing, that it behaves like answer supervision with on-policy sampling: a surface rule already earns most of the reward on a surface-learnable task (the last-digit rule is right on \ruleLastDigit\% of our prime test items), and on div7 the policy finds the procedure only if its own samples sometimes execute it \citep{yue2025rl}. We ran no RL; arm S is the nearest experiment.

\section{The compounding approximation (post hoc)}
\label{app:compound}
This derivation is post hoc: it was not part of the preregistered S5 rule, and it does not change S5's verdict. Suppose each emitted state is correct with probability $q$, independently, and that after the first error the emitted final state is uniform on $\mathbb{Z}_d$. With probability $q^k$ the trace is error-free and the answer is right. Otherwise the answer is right with probability $1/d$ on a multiple of $d$ and $(d-1)/d$ otherwise, so on a balanced test set the accuracy is $q^k+\tfrac12(1-q^k)\,[\tfrac1d+\tfrac{d-1}d]=(1+q^k)/2$. Both assumptions can fail: errors are correlated across steps (a model that miscomputes one table entry does so whenever the entry recurs), and a wrong state need not be uniform (on div13, wrong traces end at $r=0$ in \divthirteenWrongZero{} of \divthirteenWrongN{} cases). The positive-class accuracy, $q^k+(1-q^k)/d$, is the cleaner check, because a model that answers ``no'' by default scores well on non-multiples: on 6-digit div7 it is \divsevenSixByes\%.

\end{document}
